\section{Derivation concordance}
\label{app:concordance}

Every registered derivation statement, grouped by the document it lives
in, with the digest of the bytes it was read from and the heading and
line range that locate it inside that document.

\DerivStatements{} statements across \CiteDocs{} documents. The
\emph{Formal} column is \texttt{full} when a Lean theorem covers the
whole statement, \texttt{partial} when a theorem is registered as
supporting one ingredient of it, and \texttt{---} when neither. Partial
support never promotes the whole statement's tier, which is why the
two are separate columns of the underlying registry rather than one.

A caution on using this table. The status recorded here describes the
\emph{source argument} as the corpus assessed it, not an independent
re-derivation. Empty formal coverage means no complete registered
formalization exists in this inventory --- it is not a judgement that
the analytic proof is invalid, and reading it as one would invert the
meaning of the default tier (\S\ref{sec:howto}).

% Generated by paper/master/generate_sources.py.  Do not edit.
\noindent The concordance below covers all 332 registered derivation statements across 40 source documents. Each document carries the SHA-256 of the bytes its statements were read from; if the file has changed, the digest has changed and the statements must be re-located rather than assumed.\par\medskip
\subsection*{CITE:\allowbreak{}ARCHIVE\_\allowbreak{}FUNCTIONAL}
\addcontentsline{toc}{subsection}{ARCHIVE\_\allowbreak{}FUNCTIONAL}
\noindent\srcpath{docs/derivations/archive-functional-inequalities.md}\\
{\footnotesize SHA-256 \texttt{048fbf01cdc55ee4}\ldots}\par\smallskip
\begin{longtable}{@{}p{0.30\linewidth} p{0.30\linewidth} l l@{}}
\toprule Statement & Locator & Status & Formal \\ \midrule
\endfirsthead
\toprule Statement & Locator & Status & Formal \\ \midrule \endhead
\bottomrule\endfoot
{\scriptsize\ttfamily F1} & {\scriptsize \#\# F1. Trace defect, faithfulness and the vacuum set (lines 13-22)} & {\scriptsize proven} & {\scriptsize ---} \\
{\scriptsize\ttfamily F2} & {\scriptsize \#\# F2. Global gradient domination without a quotient singularity (lines 23-37)} & {\scriptsize proven} & {\scriptsize ---} \\
{\scriptsize\ttfamily F3} & {\scriptsize \#\# F3. Plaquette lifts preserve link derivatives (lines 38-53)} & {\scriptsize proven} & {\scriptsize ---} \\
{\scriptsize\ttfamily F4} & {\scriptsize \#\# F4. Aggregate defect comparisons (lines 54-64)} & {\scriptsize proven} & {\scriptsize ---} \\
{\scriptsize\ttfamily F5} & {\scriptsize \#\# F5. Exact drift identity and the squared-profile budget (lines 65-82)} & {\scriptsize proven} & {\scriptsize ---} \\
{\scriptsize\ttfamily F6} & {\scriptsize \#\# F6. General profiles with vanishing derivative at zero (lines 83-96)} & {\scriptsize proven} & {\scriptsize ---} \\
{\scriptsize\ttfamily F7} & {\scriptsize \#\# F7. Linear pairing implies a genuinely uniform small-set drift (lines 97-114)} & {\scriptsize proven} & {\scriptsize ---} \\
{\scriptsize\ttfamily F8} & {\scriptsize \#\# F8. Lyapunov energy conversion as a completed square (lines 115-129)} & {\scriptsize proven} & {\scriptsize ---} \\
{\scriptsize\ttfamily F9} & {\scriptsize \#\# F9. Local-to-global Poincare patching (lines 130-144)} & {\scriptsize proven} & {\scriptsize ---} \\
{\scriptsize\ttfamily F10} & {\scriptsize \#\# F10. Averaging gives a volume-scale gradient bound (lines 145-156)} & {\scriptsize proven} & {\scriptsize ---} \\
{\scriptsize\ttfamily F11} & {\scriptsize \#\# F11. Herbst concentration and the mean threshold (lines 157-172)} & {\scriptsize proven} & {\scriptsize ---} \\
{\scriptsize\ttfamily F12} & {\scriptsize \#\# F12. Covariance localization with an explicit error budget (lines 173-190)} & {\scriptsize proven} & {\scriptsize ---} \\
{\scriptsize\ttfamily F13} & {\scriptsize \#\# F13. Absorbing a volume tail into a distance exponent (lines 191-203)} & {\scriptsize proven} & {\scriptsize ---} \\
{\scriptsize\ttfamily F14} & {\scriptsize \#\# F14. Averaged smallness does not imply plaquettewise smallness (lines 204-213)} & {\scriptsize proven} & {\scriptsize ---} \\
{\scriptsize\ttfamily F15} & {\scriptsize \#\# F15. Center configurations refute the proposed Laplacian identity (lines 214-229)} & {\scriptsize proven} & {\scriptsize ---} \\
{\scriptsize\ttfamily F16} & {\scriptsize \#\# F16. Extensive center defects obstruct uniform pairing coercivity (lines 230-270)} & {\scriptsize proven} & {\scriptsize ---} \\
{\scriptsize\ttfamily F17} & {\scriptsize \#\# F17. The exact SU(2) diffusion budget retains favorable signs (lines 271-311)} & {\scriptsize proven} & {\scriptsize ---} \\
\end{longtable}

\subsection*{CITE:\allowbreak{}ARCHIVE\_\allowbreak{}SOURCE\_\allowbreak{}TRANSPORT}
\addcontentsline{toc}{subsection}{ARCHIVE\_\allowbreak{}SOURCE\_\allowbreak{}TRANSPORT}
\noindent\srcpath{docs/derivations/archive-source-transport.md}\\
{\footnotesize SHA-256 \texttt{4a9d953d888e1b0e}\ldots}\par\smallskip
\begin{longtable}{@{}p{0.30\linewidth} p{0.30\linewidth} l l@{}}
\toprule Statement & Locator & Status & Formal \\ \midrule
\endfirsthead
\toprule Statement & Locator & Status & Formal \\ \midrule \endhead
\bottomrule\endfoot
{\scriptsize\ttfamily S1} & {\scriptsize \#\# S1. Inhomogeneous source tilting is a partition ratio (lines 10-19)} & {\scriptsize proven} & {\scriptsize ---} \\
{\scriptsize\ttfamily S2} & {\scriptsize \#\# S2. Source response is covariance and log partition is convex (lines 20-31)} & {\scriptsize proven} & {\scriptsize ---} \\
{\scriptsize\ttfamily S3} & {\scriptsize \#\# S3. Pointwise envelopes do not order normalized tilts (lines 32-44)} & {\scriptsize proven} & {\scriptsize ---} \\
{\scriptsize\ttfamily S4} & {\scriptsize \#\# S4. Bounded log likelihood ratios transfer tilted estimates (lines 45-57)} & {\scriptsize proven} & {\scriptsize ---} \\
{\scriptsize\ttfamily S5} & {\scriptsize \#\# S5. Free energy implies conditional rooted-capacity summability (lines 58-78)} & {\scriptsize proven} & {\scriptsize ---} \\
{\scriptsize\ttfamily S6} & {\scriptsize \#\# S6. A conservative rooted-animal count (lines 79-94)} & {\scriptsize proven} & {\scriptsize ---} \\
{\scriptsize\ttfamily S7} & {\scriptsize \#\# S7. Reflection-positive deterministic pushforward (lines 95-107)} & {\scriptsize proven} & {\scriptsize ---} \\
{\scriptsize\ttfamily S8} & {\scriptsize \#\# S8. Projective and weak limits preserve the stated OS forms (lines 108-122)} & {\scriptsize proven} & {\scriptsize ---} \\
{\scriptsize\ttfamily S9} & {\scriptsize \#\# S9. Reflection equivariance alone has a counterexample (lines 123-134)} & {\scriptsize proven} & {\scriptsize ---} \\
{\scriptsize\ttfamily S10} & {\scriptsize \#\# S10. Block Schur bound with row and column control (lines 135-147)} & {\scriptsize proven} & {\scriptsize ---} \\
{\scriptsize\ttfamily S11} & {\scriptsize \#\# S11. Finite-range inverse transport with an explicit exponent (lines 148-162)} & {\scriptsize proven} & {\scriptsize ---} \\
{\scriptsize\ttfamily S12} & {\scriptsize \#\# S12. Massive Maxwell row sums from incidence counting (lines 163-178)} & {\scriptsize proven} & {\scriptsize ---} \\
{\scriptsize\ttfamily S13} & {\scriptsize \#\# S13. Cochain horizontals reduce the massive Maxwell operator (lines 179-186)} & {\scriptsize proven} & {\scriptsize ---} \\
{\scriptsize\ttfamily S14} & {\scriptsize \#\# S14. Restriction need not preserve an ambient local kernel (lines 187-201)} & {\scriptsize proven} & {\scriptsize ---} \\
\end{longtable}

\subsection*{CITE:\allowbreak{}EARLIER\_\allowbreak{}PROOF\_\allowbreak{}RECOVERY}
\addcontentsline{toc}{subsection}{EARLIER\_\allowbreak{}PROOF\_\allowbreak{}RECOVERY}
\noindent\srcpath{docs/derivations/earlier-proof-recovery.md}\\
{\footnotesize SHA-256 \texttt{4b5b067ee7697cf5}\ldots}\par\smallskip
\begin{longtable}{@{}p{0.30\linewidth} p{0.30\linewidth} l l@{}}
\toprule Statement & Locator & Status & Formal \\ \midrule
\endfirsthead
\toprule Statement & Locator & Status & Formal \\ \midrule \endhead
\bottomrule\endfoot
{\scriptsize\ttfamily CUBIC\_\allowbreak{}QUOTIENT\_\allowbreak{}REGULARITY} & {\scriptsize \#\# E1. Cubic quotient regularity trichotomy} & {\scriptsize proven} & {\scriptsize ---} \\
{\scriptsize\ttfamily EQUAL\_\allowbreak{}RAY\_\allowbreak{}IDENTIFICATION} & {\scriptsize \#\# E2. Equal-ray recovery and an independent holdout} & {\scriptsize proven} & {\scriptsize ---} \\
{\scriptsize\ttfamily LOCALIZED\_\allowbreak{}FRAME\_\allowbreak{}OBSTRUCTION} & {\scriptsize \#\# E3. No uniformly localized finite translation frame} & {\scriptsize proven} & {\scriptsize ---} \\
{\scriptsize\ttfamily JOINT\_\allowbreak{}RANK\_\allowbreak{}VOLUME\_\allowbreak{}SCALING} & {\scriptsize \#\# E4. Sharp joint rank-volume scaling} & {\scriptsize proven} & {\scriptsize ---} \\
{\scriptsize\ttfamily FINITE\_\allowbreak{}ORDER\_\allowbreak{}SUPPORT} & {\scriptsize \#\# E5a. Bounded support at a fixed order} & {\scriptsize proven} & {\scriptsize ---} \\
{\scriptsize\ttfamily GRAM\_\allowbreak{}NULL\_\allowbreak{}OPERATOR\_\allowbreak{}QUOTIENT} & {\scriptsize \#\# E5b. Gram-null decoupling} & {\scriptsize proven} & {\scriptsize ---} \\
{\scriptsize\ttfamily DECORATED\_\allowbreak{}HISTORY\_\allowbreak{}MERGING} & {\scriptsize \#\# E5c. Exact merging of decorated histories} & {\scriptsize proven} & {\scriptsize ---} \\
{\scriptsize\ttfamily NESTED\_\allowbreak{}QUOTIENT\_\allowbreak{}REDUCTION} & {\scriptsize \#\# E5d. Finite-order nested-quotient spectral reduction} & {\scriptsize proven} & {\scriptsize ---} \\
{\scriptsize\ttfamily SHELL\_\allowbreak{}SYMMETRY\_\allowbreak{}RITZ\_\allowbreak{}CONTROL} & {\scriptsize \#\# E6a. Symmetry and residual control for a reduced shell} & {\scriptsize proven} & {\scriptsize ---} \\
{\scriptsize\ttfamily B6\_\allowbreak{}RETAINED\_\allowbreak{}KRYLOV\_\allowbreak{}CAMPAIGN} & {\scriptsize \#\# E6b. Retained B6 numerical campaign} & {\scriptsize conditional} & {\scriptsize ---} \\
{\scriptsize\ttfamily ISOLATED\_\allowbreak{}POSITIVE\_\allowbreak{}MATRIX\_\allowbreak{}ATOM} & {\scriptsize \#\# E7. Isolated positive matrix atom persists} & {\scriptsize proven} & {\scriptsize ---} \\
{\scriptsize\ttfamily CONDITIONAL\_\allowbreak{}SPECTRAL\_\allowbreak{}FLOOR} & {\scriptsize \#\# E8. Conditional spectral floor and defect contraction} & {\scriptsize proven} & {\scriptsize ---} \\
{\scriptsize\ttfamily D4\_\allowbreak{}DISJOINT\_\allowbreak{}STAPLE\_\allowbreak{}COORDINATES} & {\scriptsize \#\# E9. Six disjoint staple coordinates in four dimensions} & {\scriptsize proven} & {\scriptsize ---} \\
{\scriptsize\ttfamily POSITIVE\_\allowbreak{}SECTOR\_\allowbreak{}PHASE\_\allowbreak{}ISOLATION} & {\scriptsize \#\# E10. Positive-sector phase isolation} & {\scriptsize proven} & {\scriptsize ---} \\
{\scriptsize\ttfamily VSU\_\allowbreak{}BOUNDED\_\allowbreak{}DOMAIN\_\allowbreak{}WELLPOSEDNESS} & {\scriptsize \#\# E11. Convex VSU action on a bounded domain} & {\scriptsize proven} & {\scriptsize ---} \\
{\scriptsize\ttfamily DETERMINANT\_\allowbreak{}PARITY\_\allowbreak{}REPAIR} & {\scriptsize \#\# E12. Determinant reduction with the missing parity restored} & {\scriptsize proven} & {\scriptsize ---} \\
{\scriptsize\ttfamily MOVING\_\allowbreak{}ADJOINT\_\allowbreak{}FORCE\_\allowbreak{}DERIVATIVE} & {\scriptsize \#\# E13. A moving adjoint force has two derivative terms} & {\scriptsize proven} & {\scriptsize ---} \\
{\scriptsize\ttfamily SPHERE\_\allowbreak{}EQUAL\_\allowbreak{}COUPLING\_\allowbreak{}COVARIANCE} & {\scriptsize \#\# E14. Equal-coupling SU(3) faces on a sphere have nonnegative covariance} & {\scriptsize proven} & {\scriptsize ---} \\
{\scriptsize\ttfamily SU3\_\allowbreak{}POSITIVITY\_\allowbreak{}BOUNDARIES} & {\scriptsize \#\# E15. Exact SU(3) boundaries of the covariance argument} & {\scriptsize proven} & {\scriptsize ---} \\
{\scriptsize\ttfamily REFLECTION\_\allowbreak{}ADAPTED\_\allowbreak{}BLOCKING} & {\scriptsize \#\# E16. Reflection-adapted deterministic blocking preserves positivity} & {\scriptsize proven} & {\scriptsize ---} \\
{\scriptsize\ttfamily CORNER\_\allowbreak{}BLOCKING\_\allowbreak{}REFLECTION\_\allowbreak{}DEFECT} & {\scriptsize \#\# E17. Literal corner paths fail the required reflection identity} & {\scriptsize proven} & {\scriptsize ---} \\
{\scriptsize\ttfamily DIPOLAR\_\allowbreak{}REMAINDER} & {\scriptsize \#\# E18. Explicit remaining application obligations} & {\scriptsize open} & {\scriptsize ---} \\
{\scriptsize\ttfamily PHYSICAL\_\allowbreak{}CARRIER\_\allowbreak{}REALIZATION} & {\scriptsize \#\# E18. Explicit remaining application obligations} & {\scriptsize open} & {\scriptsize ---} \\
{\scriptsize\ttfamily STAPLE\_\allowbreak{}TRANSVERSALITY\_\allowbreak{}REPAIR} & {\scriptsize \#\# E18. Explicit remaining application obligations} & {\scriptsize open} & {\scriptsize ---} \\
{\scriptsize\ttfamily VSU\_\allowbreak{}WHOLE\_\allowbreak{}SPACE} & {\scriptsize \#\# E18. Explicit remaining application obligations} & {\scriptsize open} & {\scriptsize ---} \\
{\scriptsize\ttfamily OMITTED\_\allowbreak{}DETERMINANT\_\allowbreak{}PARITY} & {\scriptsize \#\# E12. Determinant reduction with the missing parity restored} & {\scriptsize falsified} & {\scriptsize ---} \\
{\scriptsize\ttfamily OMITTED\_\allowbreak{}OWN\_\allowbreak{}FORCE\_\allowbreak{}DERIVATIVE} & {\scriptsize \#\# E13. A moving adjoint force has two derivative terms} & {\scriptsize falsified} & {\scriptsize ---} \\
\end{longtable}

\subsection*{CITE:\allowbreak{}HODGE\_\allowbreak{}POLYMER\_\allowbreak{}SUPPORTED}
\addcontentsline{toc}{subsection}{HODGE\_\allowbreak{}POLYMER\_\allowbreak{}SUPPORTED}
\noindent\srcpath{docs/derivations/hodge-polymer-supported-statements.md}\\
{\footnotesize SHA-256 \texttt{e0221c892ff185e7}\ldots}\par\smallskip
\begin{longtable}{@{}p{0.30\linewidth} p{0.30\linewidth} l l@{}}
\toprule Statement & Locator & Status & Formal \\ \midrule
\endfirsthead
\toprule Statement & Locator & Status & Formal \\ \midrule \endhead
\bottomrule\endfoot
{\scriptsize\ttfamily H1} & {\scriptsize \#\# H1: complementary kernels from Hodge relations (from line 17)} & {\scriptsize proven} & {\scriptsize full} \\
{\scriptsize\ttfamily H2} & {\scriptsize \#\# H2: carrier action and absence of Hodge coupling (from line 33)} & {\scriptsize proven} & {\scriptsize full} \\
{\scriptsize\ttfamily H3} & {\scriptsize \#\# H3: rank-one normalization, excitation and factorization (from line 45)} & {\scriptsize proven} & {\scriptsize full} \\
{\scriptsize\ttfamily H4} & {\scriptsize \#\# H4: finite word patterns and supplied polynomial identities (from line 71)} & {\scriptsize proven} & {\scriptsize full} \\
{\scriptsize\ttfamily P1} & {\scriptsize \#\# P1: geometric-series majorant (from line 93)} & {\scriptsize proven} & {\scriptsize full} \\
{\scriptsize\ttfamily P2} & {\scriptsize \#\# P2: parameter bounds and the separate tree budget (from line 107)} & {\scriptsize proven} & {\scriptsize full} \\
\end{longtable}

\subsection*{CITE:\allowbreak{}LOCAL\_\allowbreak{}CLASS\_\allowbreak{}WICK\_\allowbreak{}SPECTRUM}
\addcontentsline{toc}{subsection}{LOCAL\_\allowbreak{}CLASS\_\allowbreak{}WICK\_\allowbreak{}SPECTRUM}
\noindent\srcpath{docs/derivations/local-class-wick-spectrum.md}\\
{\footnotesize SHA-256 \texttt{5cde9ee4fe88bd47}\ldots}\par\smallskip
\begin{longtable}{@{}p{0.30\linewidth} p{0.30\linewidth} l l@{}}
\toprule Statement & Locator & Status & Formal \\ \midrule
\endfirsthead
\toprule Statement & Locator & Status & Formal \\ \midrule \endhead
\bottomrule\endfoot
{\scriptsize\ttfamily RECURRENCE} & {\scriptsize \#\# 3. All-order constructive theorem} & {\scriptsize proven} & {\scriptsize ---} \\
{\scriptsize\ttfamily COEFFICIENTS} & {\scriptsize \#\# 4. Recovered coefficients and the next odd term} & {\scriptsize proven} & {\scriptsize ---} \\
{\scriptsize\ttfamily SIGNS} & {\scriptsize \#\# 4.2 Sign corollary} & {\scriptsize proven} & {\scriptsize ---} \\
\end{longtable}

\subsection*{CITE:\allowbreak{}MOVING\_\allowbreak{}TIME\_\allowbreak{}SPECTRAL\_\allowbreak{}GAP}
\addcontentsline{toc}{subsection}{MOVING\_\allowbreak{}TIME\_\allowbreak{}SPECTRAL\_\allowbreak{}GAP}
\noindent\srcpath{docs/derivations/moving-time-spectral-gap.md}\\
{\footnotesize SHA-256 \texttt{d0a20c25e023bff7}\ldots}\par\smallskip
\begin{longtable}{@{}p{0.30\linewidth} p{0.30\linewidth} l l@{}}
\toprule Statement & Locator & Status & Formal \\ \midrule
\endfirsthead
\toprule Statement & Locator & Status & Formal \\ \midrule \endhead
\bottomrule\endfoot
{\scriptsize\ttfamily PROJECTION} & {\scriptsize \#\# 1. Premises and the finite-cutoff projection estimate (MT1) (lines 17-43)} & {\scriptsize proven} & {\scriptsize ---} \\
{\scriptsize\ttfamily VAGUE\_\allowbreak{}LIMIT} & {\scriptsize \#\# 2. Moving-time exclusion and totality (MT2) (lines 44-91)} & {\scriptsize proven} & {\scriptsize ---} \\
{\scriptsize\ttfamily SHRINKING\_\allowbreak{}TILTS} & {\scriptsize \#\# 3. Shrinking exponential sources without a Taylor rate penalty (MT3) (lines 92-126)} & {\scriptsize proven} & {\scriptsize ---} \\
{\scriptsize\ttfamily OPTIMAL\_\allowbreak{}BUDGET} & {\scriptsize \#\# 4. Optimal power-law budget and finite observation horizon (MT4) (lines 127-179)} & {\scriptsize proven} & {\scriptsize ---} \\
{\scriptsize\ttfamily SHARPNESS} & {\scriptsize \#\# 5. Sharpness and falsifiers (MT5) (lines 180-234)} & {\scriptsize proven} & {\scriptsize ---} \\
\end{longtable}

\subsection*{CITE:\allowbreak{}OS\_\allowbreak{}KERNEL\_\allowbreak{}GAP}
\addcontentsline{toc}{subsection}{OS\_\allowbreak{}KERNEL\_\allowbreak{}GAP}
\noindent\srcpath{docs/derivations/os-kernel-moving-time-gap.md}\\
{\footnotesize SHA-256 \texttt{d7088d36f67a1f8c}\ldots}\par\smallskip
\begin{longtable}{@{}p{0.30\linewidth} p{0.30\linewidth} l l@{}}
\toprule Statement & Locator & Status & Formal \\ \midrule
\endfirsthead
\toprule Statement & Locator & Status & Formal \\ \midrule \endhead
\bottomrule\endfoot
{\scriptsize\ttfamily KERNEL} & {\scriptsize \#\# K1. Construct the physical semigroup from compatible limiting kernels (lines 10-61)} & {\scriptsize proven} & {\scriptsize ---} \\
{\scriptsize\ttfamily LATE\_\allowbreak{}TIME} & {\scriptsize \#\# K2. A single late-time bound controls every limiting fixed time (lines 62-110)} & {\scriptsize proven} & {\scriptsize ---} \\
{\scriptsize\ttfamily WEIGHT} & {\scriptsize \#\# K3. Nontrivial finite-energy weight from one separated-time lower bound (lines 111-132)} & {\scriptsize proven} & {\scriptsize ---} \\
{\scriptsize\ttfamily LIMITATIONS} & {\scriptsize \#\# K4. Falsifiers and actual Wilson application check (lines 133-175)} & {\scriptsize proven} & {\scriptsize ---} \\
{\scriptsize\ttfamily POSITIVE\_\allowbreak{}HISTORY} & {\scriptsize \#\# K5. Positive-time history separation supplies zero-time continuity (lines 176-221)} & {\scriptsize proven} & {\scriptsize ---} \\
\end{longtable}

\subsection*{CITE:\allowbreak{}PBH\_\allowbreak{}WILSON\_\allowbreak{}TEST}
\addcontentsline{toc}{subsection}{PBH\_\allowbreak{}WILSON\_\allowbreak{}TEST}
\noindent\srcpath{docs/derivations/yangmills-pbh-wilson-test.md}\\
{\footnotesize SHA-256 \texttt{3122ef912820b5a0}\ldots}\par\smallskip
\begin{longtable}{@{}p{0.30\linewidth} p{0.30\linewidth} l l@{}}
\toprule Statement & Locator & Status & Formal \\ \midrule
\endfirsthead
\toprule Statement & Locator & Status & Formal \\ \midrule \endhead
\bottomrule\endfoot
{\scriptsize\ttfamily PBH1\_\allowbreak{}PBH3} & {\scriptsize \#\# PBH-1. The actual measure and horizontal form (lines 8-57)} & {\scriptsize proven} & {\scriptsize ---} \\
{\scriptsize\ttfamily PBH4} & {\scriptsize \#\# PBH-2. Explicit regular-point counterexample (lines 58-113)} & {\scriptsize proven} & {\scriptsize ---} \\
{\scriptsize\ttfamily PBH5} & {\scriptsize \#\# PBH-3. The obstruction is not confined to a one-plaquette lattice (lines 114-163)} & {\scriptsize proven} & {\scriptsize ---} \\
{\scriptsize\ttfamily PBH6} & {\scriptsize \#\# PBH-4. Classical Wilson weight is not the quantum ground-state weight (lines 164-198)} & {\scriptsize proven} & {\scriptsize ---} \\
{\scriptsize\ttfamily FLOW\_\allowbreak{}COUNTEREXAMPLE} & {\scriptsize \#\# PBH-5. The stated gradient-flow shortcut also fails an exact test (lines 199-230)} & {\scriptsize proven} & {\scriptsize ---} \\
\end{longtable}

\subsection*{CITE:\allowbreak{}REVIEW\_\allowbreak{}SPECTRAL\_\allowbreak{}BUDGETS}
\addcontentsline{toc}{subsection}{REVIEW\_\allowbreak{}SPECTRAL\_\allowbreak{}BUDGETS}
\noindent\srcpath{docs/derivations/review-derived-spectral-budgets.md}\\
{\footnotesize SHA-256 \texttt{96037da9bfc02943}\ldots}\par\smallskip
\begin{longtable}{@{}p{0.30\linewidth} p{0.30\linewidth} l l@{}}
\toprule Statement & Locator & Status & Formal \\ \midrule
\endfirsthead
\toprule Statement & Locator & Status & Formal \\ \midrule \endhead
\bottomrule\endfoot
{\scriptsize\ttfamily Q1} & {\scriptsize \#\# Q1. The assembled fourth-order polynomial is coercive (lines 11-42)} & {\scriptsize proven} & {\scriptsize ---} \\
{\scriptsize\ttfamily M1} & {\scriptsize \#\# M1. Several decay channels and time-amplified remainders (lines 43-89)} & {\scriptsize proven} & {\scriptsize ---} \\
{\scriptsize\ttfamily M2} & {\scriptsize \#\# M2. Closed optimal rate and a witnessing observation time (lines 90-125)} & {\scriptsize proven} & {\scriptsize ---} \\
{\scriptsize\ttfamily M3} & {\scriptsize \#\# M3. Sharpness, failure cases and physical scope (lines 126-171)} & {\scriptsize proven} & {\scriptsize ---} \\
\end{longtable}

\subsection*{CITE:\allowbreak{}SEPTEMBER10\_\allowbreak{}UPSTREAM\_\allowbreak{}SCALARS}
\addcontentsline{toc}{subsection}{SEPTEMBER10\_\allowbreak{}UPSTREAM\_\allowbreak{}SCALARS}
\noindent\srcpath{docs/derivations/september10-upstream-scalar-inventory.md}\\
{\footnotesize SHA-256 \texttt{75a6b78ebaed289f}\ldots}\par\smallskip
\begin{longtable}{@{}p{0.30\linewidth} p{0.30\linewidth} l l@{}}
\toprule Statement & Locator & Status & Formal \\ \midrule
\endfirsthead
\toprule Statement & Locator & Status & Formal \\ \midrule \endhead
\bottomrule\endfoot
{\scriptsize\ttfamily KP\_\allowbreak{}SCALARS} & {\scriptsize \#\# U1 Tilted KP scalar ingredients} & {\scriptsize proven} & {\scriptsize partial} \\
{\scriptsize\ttfamily COERCIVITY\_\allowbreak{}SCALARS} & {\scriptsize \#\# U2 Coercivity scalar ingredients} & {\scriptsize proven} & {\scriptsize partial} \\
\end{longtable}

\subsection*{CITE:\allowbreak{}SEPT\_\allowbreak{}DOC\_\allowbreak{}YANGMILLS\_\allowbreak{}FORMAL\_\allowbreak{}BRIDGES}
\addcontentsline{toc}{subsection}{SEPT\_\allowbreak{}DOC\_\allowbreak{}YANGMILLS\_\allowbreak{}FORMAL\_\allowbreak{}BRIDGES}
\noindent\srcpath{docs/derivations/yangmills-formal-bridges.md}\\
{\footnotesize SHA-256 \texttt{bb3245e82e5dc857}\ldots}\par\smallskip
\begin{longtable}{@{}p{0.30\linewidth} p{0.30\linewidth} l l@{}}
\toprule Statement & Locator & Status & Formal \\ \midrule
\endfirsthead
\toprule Statement & Locator & Status & Formal \\ \midrule \endhead
\bottomrule\endfoot
{\scriptsize\ttfamily F1} & {\scriptsize \#\# Exact finite ground-state identification (lines 14-78)} & {\scriptsize proven} & {\scriptsize full} \\
{\scriptsize\ttfamily F2} & {\scriptsize \#\# Exact finite ground-state identification (lines 14-78)} & {\scriptsize proven} & {\scriptsize full} \\
{\scriptsize\ttfamily F3} & {\scriptsize \#\# Exact finite ground-state identification (lines 14-78)} & {\scriptsize proven} & {\scriptsize full} \\
{\scriptsize\ttfamily F4} & {\scriptsize \#\# A trial residual needs a floor and a mean (lines 79-101)} & {\scriptsize proven} & {\scriptsize full} \\
{\scriptsize\ttfamily F5} & {\scriptsize \#\# The pure quartic Gaussian certificate (lines 102-137)} & {\scriptsize proven} & {\scriptsize full} \\
{\scriptsize\ttfamily F6} & {\scriptsize \#\# The pure quartic Gaussian certificate (lines 102-137)} & {\scriptsize proven} & {\scriptsize full} \\
{\scriptsize\ttfamily F7} & {\scriptsize \#\# Algebra for Simon's transverse mechanism (lines 138-169)} & {\scriptsize proven} & {\scriptsize full} \\
{\scriptsize\ttfamily SLICE\_\allowbreak{}RADIAL\_\allowbreak{}ALGEBRA} & {\scriptsize \#\# Algebra for Simon's transverse mechanism (lines 138-169)} & {\scriptsize proven} & {\scriptsize ---} \\
\end{longtable}

\subsection*{CITE:\allowbreak{}SEPT\_\allowbreak{}DOC\_\allowbreak{}YANGMILLS\_\allowbreak{}RESOLVENT\_\allowbreak{}LOCALIZATION\_\allowbreak{}AND\_\allowbreak{}DIRICHLET\_\allowbreak{}GAP}
\addcontentsline{toc}{subsection}{SEPT\_\allowbreak{}DOC\_\allowbreak{}YANGMILLS\_\allowbreak{}RESOLVENT\_\allowbreak{}LOCALIZATION\_\allowbreak{}AND\_\allowbreak{}DIRICHLET\_\allowbreak{}GAP}
\noindent\srcpath{docs/derivations/yangmills-resolvent-localization-and-dirichlet-gap.md}\\
{\footnotesize SHA-256 \texttt{876b76c1aee3eb8e}\ldots}\par\smallskip
\begin{longtable}{@{}p{0.30\linewidth} p{0.30\linewidth} l l@{}}
\toprule Statement & Locator & Status & Formal \\ \midrule
\endfirsthead
\toprule Statement & Locator & Status & Formal \\ \midrule \endhead
\bottomrule\endfoot
{\scriptsize\ttfamily SUBMITTED\_\allowbreak{}W6\_\allowbreak{}CLOSURE} & {\scriptsize \#\# 1. Summary of Results (lines 20-39)} & {\scriptsize disputed} & {\scriptsize ---} \\
{\scriptsize\ttfamily RESOLVENT\_\allowbreak{}IDENTITY} & {\scriptsize \#\#\# 2.1 The Resolvent Identity (lines 42-54)} & {\scriptsize conditional} & {\scriptsize partial} \\
{\scriptsize\ttfamily LOCAL\_\allowbreak{}KERNEL\_\allowbreak{}IMPLICATION} & {\scriptsize \#\#\# 2.3 Volume-Uniformity of the Pairing (lines 62-80)} & {\scriptsize conditional} & {\scriptsize partial} \\
{\scriptsize\ttfamily GROUND\_\allowbreak{}TRANSFORM} & {\scriptsize \#\#\# 3.2 True Ground Conjugation and Exact Cancellation (lines 104-123)} & {\scriptsize proven} & {\scriptsize ---} \\
{\scriptsize\ttfamily SUBMITTED\_\allowbreak{}CLOSURES} & {\scriptsize \#\# 4. Status and Graph Integration (lines 131-135)} & {\scriptsize superseded} & {\scriptsize ---} \\
\end{longtable}

\subsection*{CITE:\allowbreak{}SOURCE\_\allowbreak{}CURRENT\_\allowbreak{}SPECTRAL\_\allowbreak{}BRIDGES}
\addcontentsline{toc}{subsection}{SOURCE\_\allowbreak{}CURRENT\_\allowbreak{}SPECTRAL\_\allowbreak{}BRIDGES}
\noindent\srcpath{docs/derivations/source-currents-and-spectral-totality.md}\\
{\footnotesize SHA-256 \texttt{860c81a3fb665d6b}\ldots}\par\smallskip
\begin{longtable}{@{}p{0.30\linewidth} p{0.30\linewidth} l l@{}}
\toprule Statement & Locator & Status & Formal \\ \midrule
\endfirsthead
\toprule Statement & Locator & Status & Formal \\ \midrule \endhead
\bottomrule\endfoot
{\scriptsize\ttfamily SCB0A} & {\scriptsize \#\# SCB0A. Degenerate energy variational certificate (lines 21-34)} & {\scriptsize proven} & {\scriptsize partial} \\
{\scriptsize\ttfamily SCB0B} & {\scriptsize \#\# SCB0B. Bounded-inverse W6 composition (lines 35-59)} & {\scriptsize proven} & {\scriptsize partial} \\
{\scriptsize\ttfamily SCB1} & {\scriptsize \#\# SCB1. The complete first current in the actual square convention (lines 60-114)} & {\scriptsize proven} & {\scriptsize ---} \\
{\scriptsize\ttfamily SCB2} & {\scriptsize \#\# SCB2. Exact current moments and an all-source constant (lines 115-172)} & {\scriptsize proven} & {\scriptsize partial} \\
{\scriptsize\ttfamily SCB3} & {\scriptsize \#\# SCB3. A selected-inverse implication without a separate diagonal bound (lines 173-197)} & {\scriptsize proven} & {\scriptsize partial} \\
{\scriptsize\ttfamily SCB4} & {\scriptsize \#\# SCB4. Intrinsic score currents and exact centering corrections (lines 198-257)} & {\scriptsize proven} & {\scriptsize ---} \\
{\scriptsize\ttfamily SCB5} & {\scriptsize \#\# SCB5. Tiny exponential sources imply totality by differentiation (lines 258-298)} & {\scriptsize proven} & {\scriptsize partial} \\
{\scriptsize\ttfamily SCB6} & {\scriptsize \#\# SCB6. Uniform normalized mesoscopic source amplitudes (lines 299-316)} & {\scriptsize proven} & {\scriptsize ---} \\
\end{longtable}

\subsection*{CITE:\allowbreak{}THEORY\_\allowbreak{}GEOMETRY\_\allowbreak{}RICCATI\_\allowbreak{}BRIDGES}
\addcontentsline{toc}{subsection}{THEORY\_\allowbreak{}GEOMETRY\_\allowbreak{}RICCATI\_\allowbreak{}BRIDGES}
\noindent\srcpath{docs/derivations/theory-geometry-and-riccati-bridges.md}\\
{\footnotesize SHA-256 \texttt{7216dd8c2f263782}\ldots}\par\smallskip
\begin{longtable}{@{}p{0.30\linewidth} p{0.30\linewidth} l l@{}}
\toprule Statement & Locator & Status & Formal \\ \midrule
\endfirsthead
\toprule Statement & Locator & Status & Formal \\ \midrule \endhead
\bottomrule\endfoot
{\scriptsize\ttfamily TG1} & {\scriptsize \#\# TG1: The sharp three-coordinate anisotropy constant (lines 14-60)} & {\scriptsize proven} & {\scriptsize partial} \\
{\scriptsize\ttfamily TG2} & {\scriptsize \#\# TG2: Zero extension and integrated energy bounds (lines 61-93)} & {\scriptsize proven} & {\scriptsize ---} \\
{\scriptsize\ttfamily TG3} & {\scriptsize \#\# TG3: One secular cubic for the recorded Hodge support (lines 94-136)} & {\scriptsize proven} & {\scriptsize partial} \\
{\scriptsize\ttfamily TG4} & {\scriptsize \#\# TG4: The induced eighth-order moment in the same truncation (lines 137-166)} & {\scriptsize proven} & {\scriptsize partial} \\
{\scriptsize\ttfamily TG5} & {\scriptsize \#\# TG5: Complete source motion is a Schur congruence (lines 167-222)} & {\scriptsize proven} & {\scriptsize partial} \\
{\scriptsize\ttfamily TG6} & {\scriptsize \#\# TG6: The metric and finite source resolvent must travel with the form (lines 223-274)} & {\scriptsize proven} & {\scriptsize partial} \\
{\scriptsize\ttfamily TG7} & {\scriptsize \#\# TG7: Scalar square-source parity eliminates formal odd Schur jets (lines 275-313)} & {\scriptsize proven} & {\scriptsize ---} \\
{\scriptsize\ttfamily TG8} & {\scriptsize \#\# TG8: A residual certificate for the abstract Riccati fixed point (lines 314-350)} & {\scriptsize proven} & {\scriptsize partial} \\
\end{longtable}

\subsection*{CITE:\allowbreak{}TRACK\_\allowbreak{}A\_\allowbreak{}SUPPORTED}
\addcontentsline{toc}{subsection}{TRACK\_\allowbreak{}A\_\allowbreak{}SUPPORTED}
\noindent\srcpath{docs/derivations/track-a-supported-statements.md}\\
{\footnotesize SHA-256 \texttt{26efd60d345eb4c9}\ldots}\par\smallskip
\begin{longtable}{@{}p{0.30\linewidth} p{0.30\linewidth} l l@{}}
\toprule Statement & Locator & Status & Formal \\ \midrule
\endfirsthead
\toprule Statement & Locator & Status & Formal \\ \midrule \endhead
\bottomrule\endfoot
{\scriptsize\ttfamily W6\_\allowbreak{}V} & {\scriptsize \#\# W6-V: actual variational residual identity and dual energy (from line 18)} & {\scriptsize proven} & {\scriptsize full} \\
{\scriptsize\ttfamily W6\_\allowbreak{}F} & {\scriptsize \#\# W6-F: bounded-operator factorization controls the residual functional (from line 43)} & {\scriptsize proven} & {\scriptsize full} \\
{\scriptsize\ttfamily W6\_\allowbreak{}E} & {\scriptsize \#\# W6-E: sufficient selected-inverse error certificate (from line 61)} & {\scriptsize proven} & {\scriptsize full} \\
{\scriptsize\ttfamily W6\_\allowbreak{}C} & {\scriptsize \#\# W6-C: exact Gaussian coefficient substitutions (from line 76)} & {\scriptsize proven} & {\scriptsize full} \\
{\scriptsize\ttfamily SC17\_\allowbreak{}R} & {\scriptsize \#\# SC17-R: Riccati roots and the continuous barrier (from line 92)} & {\scriptsize proven} & {\scriptsize full} \\
{\scriptsize\ttfamily SC17\_\allowbreak{}B} & {\scriptsize \#\# SC17-B: constructed noncommutative Banach-algebra fixed point (from line 112)} & {\scriptsize proven} & {\scriptsize full} \\
{\scriptsize\ttfamily SC17\_\allowbreak{}P} & {\scriptsize \#\# SC17-P: exact interval arithmetic under the proposed inputs (from line 132)} & {\scriptsize proven} & {\scriptsize full} \\
{\scriptsize\ttfamily G17\_\allowbreak{}P} & {\scriptsize \#\# G17-P: corrected raw partition estimate for actual probability measures (from line 151)} & {\scriptsize proven} & {\scriptsize full} \\
{\scriptsize\ttfamily G17\_\allowbreak{}V} & {\scriptsize \#\# G17-V: centered L2 source, variance, and bounded footprints (from line 168)} & {\scriptsize proven} & {\scriptsize full} \\
{\scriptsize\ttfamily G17\_\allowbreak{}O} & {\scriptsize \#\# G17-O: strict obstruction to a constant equal to one (from line 185)} & {\scriptsize proven} & {\scriptsize full} \\
{\scriptsize\ttfamily G17\_\allowbreak{}I} & {\scriptsize \#\# G17-I: independent product sources have unbounded radius (from line 206)} & {\scriptsize proven} & {\scriptsize full} \\
\end{longtable}

\subsection*{CITE:\allowbreak{}UNIVERSAL\_\allowbreak{}CELLULAR\_\allowbreak{}HODGE\_\allowbreak{}TETRAHEDRAL}
\addcontentsline{toc}{subsection}{UNIVERSAL\_\allowbreak{}CELLULAR\_\allowbreak{}HODGE\_\allowbreak{}TETRAHEDRAL}
\noindent\srcpath{docs/derivations/universal-cellular-hodge-tetrahedral.md}\\
{\footnotesize SHA-256 \texttt{427974f8eda3ae77}\ldots}\par\smallskip
\begin{longtable}{@{}p{0.30\linewidth} p{0.30\linewidth} l l@{}}
\toprule Statement & Locator & Status & Formal \\ \midrule
\endfirsthead
\toprule Statement & Locator & Status & Formal \\ \midrule \endhead
\bottomrule\endfoot
{\scriptsize\ttfamily SPECTRUM\_\allowbreak{}EIGENVALUE} & {\scriptsize \#\#\# 2.2 Universal spectral identities (lines 45-58)} & {\scriptsize proven} & {\scriptsize partial} \\
{\scriptsize\ttfamily UP\_\allowbreak{}HARMONIC\_\allowbreak{}EXCURSION} & {\scriptsize \#\#\# 2.3 Universal Up-Harmonicity of Excursions (lines 59-67)} & {\scriptsize proven} & {\scriptsize full} \\
{\scriptsize\ttfamily TOTAL\_\allowbreak{}LAPLACIAN\_\allowbreak{}DIAGONAL} & {\scriptsize \#\#\# 2.4 Total Laplacian Diagonal \& Regularity Criterion (lines 68-83)} & {\scriptsize proven} & {\scriptsize ---} \\
{\scriptsize\ttfamily TETRAHEDRAL\_\allowbreak{}S4\_\allowbreak{}COMMUTANT} & {\scriptsize \#\#\# 3.1 Tetrahedral Hodge duality and 3.2 The S4 commutant (lines 86-119)} & {\scriptsize proven} & {\scriptsize partial} \\
{\scriptsize\ttfamily PROPER\_\allowbreak{}RETURN\_\allowbreak{}Q\_\allowbreak{}ANNIHILATION} & {\scriptsize \#\#\# 3.4 Open physical-history identification (lines 133-156)} & {\scriptsize open} & {\scriptsize partial} \\
{\scriptsize\ttfamily RETAINED\_\allowbreak{}VECTOR\_\allowbreak{}RANK\_\allowbreak{}ONE} & {\scriptsize \#\#\# 3.3 Feshbach Intermediate Annihilation of Retained Vectors (lines 120-132)} & {\scriptsize proven} & {\scriptsize full} \\
{\scriptsize\ttfamily RSMR\_\allowbreak{}MASTER\_\allowbreak{}THEOREM} & {\scriptsize \#\#\# 4.3 The carrier symbol (lines 196-221)} & {\scriptsize proven} & {\scriptsize partial} \\
{\scriptsize\ttfamily SHIFTED\_\allowbreak{}OPERATOR\_\allowbreak{}POWERS} & {\scriptsize \#\#\# 4.2 Polynomial recurrence and operator powers (lines 172-195)} & {\scriptsize proven} & {\scriptsize full} \\
\end{longtable}

\subsection*{CITE:\allowbreak{}W6\_\allowbreak{}ANTIPODAL\_\allowbreak{}MAGNETIC\_\allowbreak{}GEOMETRY}
\addcontentsline{toc}{subsection}{W6\_\allowbreak{}ANTIPODAL\_\allowbreak{}MAGNETIC\_\allowbreak{}GEOMETRY}
\noindent\srcpath{docs/derivations/w6-antipodal-magnetic-geometry.md}\\
{\footnotesize SHA-256 \texttt{5b3e0b2312f11ee8}\ldots}\par\smallskip
\begin{longtable}{@{}p{0.30\linewidth} p{0.30\linewidth} l l@{}}
\toprule Statement & Locator & Status & Formal \\ \midrule
\endfirsthead
\toprule Statement & Locator & Status & Formal \\ \midrule \endhead
\bottomrule\endfoot
{\scriptsize\ttfamily HESSIAN\_\allowbreak{}SPECTRUM} & {\scriptsize A1-A6; exact fixed-Q chart, quadratic form and full spectrum} & {\scriptsize proven} & {\scriptsize ---} \\
{\scriptsize\ttfamily GAUGE\_\allowbreak{}NORMAL\_\allowbreak{}COERCIVITY} & {\scriptsize A7-A11; complete minimum sphere, tangent kernel and original electric metric} & {\scriptsize proven} & {\scriptsize ---} \\
{\scriptsize\ttfamily ANGULAR\_\allowbreak{}REDUCTION} & {\scriptsize A12-A13; exact old-orbit restriction and compact normal-bundle implicit function} & {\scriptsize proven} & {\scriptsize ---} \\
{\scriptsize\ttfamily GAUGE\_\allowbreak{}SCORE\_\allowbreak{}INVARIANCE} & {\scriptsize A14-A15; actual ground and Haar half-density symmetry} & {\scriptsize proven} & {\scriptsize ---} \\
\end{longtable}

\subsection*{CITE:\allowbreak{}W6\_\allowbreak{}CONDITIONAL\_\allowbreak{}SCORE\_\allowbreak{}TAIL\_\allowbreak{}CONTROL}
\addcontentsline{toc}{subsection}{W6\_\allowbreak{}CONDITIONAL\_\allowbreak{}SCORE\_\allowbreak{}TAIL\_\allowbreak{}CONTROL}
\noindent\srcpath{docs/derivations/w6-conditional-score-tail-control.md}\\
{\footnotesize SHA-256 \texttt{d04ae01d29d33cc5}\ldots}\par\smallskip
\begin{longtable}{@{}p{0.30\linewidth} p{0.30\linewidth} l l@{}}
\toprule Statement & Locator & Status & Formal \\ \midrule
\endfirsthead
\toprule Statement & Locator & Status & Formal \\ \midrule \endhead
\bottomrule\endfoot
{\scriptsize\ttfamily INTERVAL\_\allowbreak{}OBSTRUCTION\_\allowbreak{}REPAIR} & {\scriptsize \#\# 1. What the actual compact endpoints do prove (lines 67-284)} & {\scriptsize proven} & {\scriptsize ---} \\
{\scriptsize\ttfamily SCORE\_\allowbreak{}DOMINATION\_\allowbreak{}M10} & {\scriptsize \#\#\# M.4. The precise remaining conditional-score inequality (lines 473-495)} & {\scriptsize open} & {\scriptsize ---} \\
{\scriptsize\ttfamily ACTUAL\_\allowbreak{}SOURCE\_\allowbreak{}MOMENT} & {\scriptsize \#\# 6. Uniform actual source potential moment and sharper score criterion (lines 325-586)} & {\scriptsize proven} & {\scriptsize ---} \\
{\scriptsize\ttfamily SCORE\_\allowbreak{}WEIGHT\_\allowbreak{}FROM\_\allowbreak{}M10} & {\scriptsize \#\#\# M.4. The precise remaining conditional-score inequality and M.5 (lines 497-563)} & {\scriptsize conditional} & {\scriptsize ---} \\
\end{longtable}

\subsection*{CITE:\allowbreak{}W6\_\allowbreak{}CONDITIONAL\_\allowbreak{}TRANSPORT\_\allowbreak{}OBSTRUCTION}
\addcontentsline{toc}{subsection}{W6\_\allowbreak{}CONDITIONAL\_\allowbreak{}TRANSPORT\_\allowbreak{}OBSTRUCTION}
\noindent\srcpath{docs/derivations/w6-conditional-transport-obstruction.md}\\
{\footnotesize SHA-256 \texttt{33f5cefc2105ec3e}\ldots}\par\smallskip
\begin{longtable}{@{}p{0.30\linewidth} p{0.30\linewidth} l l@{}}
\toprule Statement & Locator & Status & Formal \\ \midrule
\endfirsthead
\toprule Statement & Locator & Status & Formal \\ \midrule \endhead
\bottomrule\endfoot
{\scriptsize\ttfamily SCORE\_\allowbreak{}CONTINUITY} & {\scriptsize \#\# S1 Exact conditional score (lines 15-72)} & {\scriptsize proven} & {\scriptsize ---} \\
{\scriptsize\ttfamily AGMON\_\allowbreak{}GEOMETRY} & {\scriptsize \#\# S2 Exact constrained geometry (lines 73-140)} & {\scriptsize proven} & {\scriptsize ---} \\
{\scriptsize\ttfamily CONDITIONAL\_\allowbreak{}CONCENTRATION} & {\scriptsize \#\# S3 Actual conditional concentration (lines 141-213)} & {\scriptsize proven} & {\scriptsize ---} \\
{\scriptsize\ttfamily ARBITRARY\_\allowbreak{}Q8\_\allowbreak{}OBSTRUCTION} & {\scriptsize \#\# S4 A transport obstruction (lines 214-284)} & {\scriptsize proven} & {\scriptsize ---} \\
{\scriptsize\ttfamily SYNCHRONIZED\_\allowbreak{}TANGENCY} & {\scriptsize \#\# S5 Synchronized tangency (lines 285-316)} & {\scriptsize proven} & {\scriptsize ---} \\
{\scriptsize\ttfamily LOCAL\_\allowbreak{}GLOBAL\_\allowbreak{}SCORE\_\allowbreak{}CRITERION} & {\scriptsize \#\# S6 Remaining all-fiber estimate (lines 317-375)} & {\scriptsize conditional} & {\scriptsize ---} \\
{\scriptsize\ttfamily M10\_\allowbreak{}SYNCHRONIZED\_\allowbreak{}SCORE} & {\scriptsize \#\# S6 Remaining all-fiber estimate (lines 317-375), explicitly open synchronized M10 target} & {\scriptsize open} & {\scriptsize ---} \\
\end{longtable}

\subsection*{CITE:\allowbreak{}W6\_\allowbreak{}GROUND\_\allowbreak{}JETS\_\allowbreak{}TRANSPORT\_\allowbreak{}BUDGET}
\addcontentsline{toc}{subsection}{W6\_\allowbreak{}GROUND\_\allowbreak{}JETS\_\allowbreak{}TRANSPORT\_\allowbreak{}BUDGET}
\noindent\srcpath{docs/derivations/w6-ground-jets-and-transport-budget.md}\\
{\footnotesize SHA-256 \texttt{a1638b09ba941b0c}\ldots}\par\smallskip
\begin{longtable}{@{}p{0.30\linewidth} p{0.30\linewidth} l l@{}}
\toprule Statement & Locator & Status & Formal \\ \midrule
\endfirsthead
\toprule Statement & Locator & Status & Formal \\ \midrule \endhead
\bottomrule\endfoot
{\scriptsize\ttfamily GROUND\_\allowbreak{}JETS} & {\scriptsize \#\# 2. Raw electric and magnetic derivatives already have the required scaling (lines 51-198)} & {\scriptsize proven} & {\scriptsize ---} \\
{\scriptsize\ttfamily SOURCE\_\allowbreak{}ENERGY\_\allowbreak{}JETS\_\allowbreak{}R10} & {\scriptsize \#\# 4. An explicit bridge from the actual transport generator to M1,M2,M3 (lines 199-218); section 5 (lines 271-309)} & {\scriptsize open} & {\scriptsize ---} \\
{\scriptsize\ttfamily COMPLETE\_\allowbreak{}BUDGET} & {\scriptsize \#\# 4. An explicit bridge from the actual transport generator to M1,M2,M3 (lines 199-367)} & {\scriptsize conditional} & {\scriptsize ---} \\
{\scriptsize\ttfamily INTERACTING\_\allowbreak{}GRID\_\allowbreak{}COMPARISON} & {\scriptsize \#\# 6. Source energies, subdivision, interacting grids, and the physical clock (lines 351-366)} & {\scriptsize open} & {\scriptsize ---} \\
\end{longtable}

\subsection*{CITE:\allowbreak{}W6\_\allowbreak{}SOURCE\_\allowbreak{}ENERGY\_\allowbreak{}JETS\_\allowbreak{}R10}
\addcontentsline{toc}{subsection}{W6\_\allowbreak{}SOURCE\_\allowbreak{}ENERGY\_\allowbreak{}JETS\_\allowbreak{}R10}
\noindent\srcpath{docs/derivations/w6-source-energy-jets-r10.md}\\
{\footnotesize SHA-256 \texttt{7cb2ad710a35e1dd}\ldots}\par\smallskip
\begin{longtable}{@{}p{0.30\linewidth} p{0.30\linewidth} l l@{}}
\toprule Statement & Locator & Status & Formal \\ \midrule
\endfirsthead
\toprule Statement & Locator & Status & Formal \\ \midrule \endhead
\bottomrule\endfoot
{\scriptsize\ttfamily H0\_\allowbreak{}ENERGY\_\allowbreak{}CONTRACTION} & {\scriptsize \#\# 6. Scaling defects and unresolved estimates (C6)} & {\scriptsize open} & {\scriptsize ---} \\
{\scriptsize\ttfamily H1\_\allowbreak{}FIBERWISE\_\allowbreak{}SCORE\_\allowbreak{}VARIANCE} & {\scriptsize \#\# 6. Scaling defects and unresolved estimates (C6)} & {\scriptsize open} & {\scriptsize ---} \\
{\scriptsize\ttfamily VACUUM\_\allowbreak{}CROSS\_\allowbreak{}DERIVATIVES} & {\scriptsize \#\# 6. Scaling defects and unresolved estimates (C6)} & {\scriptsize open} & {\scriptsize ---} \\
{\scriptsize\ttfamily KATO\_\allowbreak{}DERIVATIVES} & {\scriptsize \#\# 6. Scaling defects and unresolved estimates (C6)} & {\scriptsize open} & {\scriptsize ---} \\
{\scriptsize\ttfamily CONDITIONAL\_\allowbreak{}DERIVATIVE} & {\scriptsize \#\# 2. Correct conditional differentiation (C1)} & {\scriptsize proven} & {\scriptsize ---} \\
{\scriptsize\ttfamily HORIZONTAL\_\allowbreak{}SCORE\_\allowbreak{}SUFFICIENT} & {\scriptsize \#\# 3. A sufficient horizontal-score estimate for H0 (C3)} & {\scriptsize proven} & {\scriptsize ---} \\
{\scriptsize\ttfamily H4\_\allowbreak{}FROM\_\allowbreak{}H0} & {\scriptsize \#\# 4. H4 follows from H0 and the established ground jet (C4)} & {\scriptsize proven} & {\scriptsize ---} \\
{\scriptsize\ttfamily PROJECTION\_\allowbreak{}JET\_\allowbreak{}RECURRENCE} & {\scriptsize \#\# 5. Complete projection-jet recurrence (C5)} & {\scriptsize proven} & {\scriptsize ---} \\
{\scriptsize\ttfamily SCALING\_\allowbreak{}OBSTRUCTIONS} & {\scriptsize \#\# 6. Scaling defects and unresolved estimates (C6)} & {\scriptsize proven} & {\scriptsize ---} \\
\end{longtable}

\subsection*{CITE:\allowbreak{}W6\_\allowbreak{}SOURCE\_\allowbreak{}GENERATOR\_\allowbreak{}SCORE\_\allowbreak{}FRAME}
\addcontentsline{toc}{subsection}{W6\_\allowbreak{}SOURCE\_\allowbreak{}GENERATOR\_\allowbreak{}SCORE\_\allowbreak{}FRAME}
\noindent\srcpath{docs/derivations/w6-source-generator-score-frame.md}\\
{\footnotesize SHA-256 \texttt{20b070cb78311187}\ldots}\par\smallskip
\begin{longtable}{@{}p{0.30\linewidth} p{0.30\linewidth} l l@{}}
\toprule Statement & Locator & Status & Formal \\ \midrule
\endfirsthead
\toprule Statement & Locator & Status & Formal \\ \midrule \endhead
\bottomrule\endfoot
{\scriptsize\ttfamily GROUND\_\allowbreak{}FRAME\_\allowbreak{}GENERATOR} & {\scriptsize \#\# 2. The ground-state frame (SF1-SF2) and \#\# 3. The complete R9 generator in the ground-state frame (SF3-SF5)} & {\scriptsize proven} & {\scriptsize ---} \\
{\scriptsize\ttfamily SCORE\_\allowbreak{}POISSON} & {\scriptsize \#\# 4. The score Poisson equation (SF6-SF7)} & {\scriptsize proven} & {\scriptsize ---} \\
{\scriptsize\ttfamily FIBER\_\allowbreak{}ENERGY\_\allowbreak{}IDENTITY} & {\scriptsize \#\# 4. The score Poisson equation (SF6-SF7), conditional divergence identity and fiber identity} & {\scriptsize proven} & {\scriptsize ---} \\
{\scriptsize\ttfamily R10\_\allowbreak{}ORDER\_\allowbreak{}ZERO\_\allowbreak{}REDUCTION} & {\scriptsize \#\# 5. Exact reduction of the order-zero R10 bound (SF8-SF9)} & {\scriptsize conditional} & {\scriptsize ---} \\
{\scriptsize\ttfamily CONDITIONAL\_\allowbreak{}ENERGY\_\allowbreak{}H0\_\allowbreak{}H4} & {\scriptsize \#\# 5. Exact reduction of the order-zero R10 bound (SF8-SF9), conditional-energy hypotheses} & {\scriptsize open} & {\scriptsize ---} \\
\end{longtable}

\subsection*{CITE:\allowbreak{}W6\_\allowbreak{}SYNCHRONIZED\_\allowbreak{}M10\_\allowbreak{}DOMINATION}
\addcontentsline{toc}{subsection}{W6\_\allowbreak{}SYNCHRONIZED\_\allowbreak{}M10\_\allowbreak{}DOMINATION}
\noindent\srcpath{docs/derivations/w6-synchronized-m10-domination.md}\\
{\footnotesize SHA-256 \texttt{e81031e9f7d13467}\ldots}\par\smallskip
\begin{longtable}{@{}p{0.30\linewidth} p{0.30\linewidth} l l@{}}
\toprule Statement & Locator & Status & Formal \\ \midrule
\endfirsthead
\toprule Statement & Locator & Status & Formal \\ \midrule \endhead
\bottomrule\endfoot
{\scriptsize\ttfamily TANGENCY\_\allowbreak{}DRIFT\_\allowbreak{}CANCELLATION} & {\scriptsize \#\# 2. Synchronized tangency (lines 34-51)} & {\scriptsize proven} & {\scriptsize ---} \\
{\scriptsize\ttfamily TUBE\_\allowbreak{}VARIANCE\_\allowbreak{}DOMINATION} & {\scriptsize \#\# 3. Small-angle tube implication (lines 52-84)} & {\scriptsize proven} & {\scriptsize ---} \\
{\scriptsize\ttfamily ANTIPODAL\_\allowbreak{}SCORE\_\allowbreak{}REGULARITY} & {\scriptsize \#\# 4. Antipodal geometry and a uniform angular reference calculation (lines 85-143)} & {\scriptsize proven} & {\scriptsize ---} \\
{\scriptsize\ttfamily AMPLITUDE\_\allowbreak{}GRADIENT\_\allowbreak{}SCALING} & {\scriptsize \#\# 5. Correct parameter-derivative criterion (lines 144-171)} & {\scriptsize open} & {\scriptsize ---} \\
{\scriptsize\ttfamily OUTSIDE\_\allowbreak{}COMPLEMENT\_\allowbreak{}SUPPRESSION} & {\scriptsize \#\# 6. Correct complement argument and compact action gap (lines 172-261)} & {\scriptsize proven} & {\scriptsize ---} \\
{\scriptsize\ttfamily FULL\_\allowbreak{}SCORE\_\allowbreak{}DOMINATION\_\allowbreak{}M10} & {\scriptsize \#\# 7. Assembly criterion and actual status (lines 262-286)} & {\scriptsize open} & {\scriptsize ---} \\
{\scriptsize\ttfamily LOG\_\allowbreak{}PARAMETER\_\allowbreak{}CRITERION} & {\scriptsize \#\# 5. Correct parameter-derivative criterion (lines 144-171)} & {\scriptsize proven} & {\scriptsize ---} \\
{\scriptsize\ttfamily COMPACT\_\allowbreak{}ACTION\_\allowbreak{}GAP} & {\scriptsize \#\# 6. Correct complement argument and compact action gap (lines 172-261)} & {\scriptsize proven} & {\scriptsize ---} \\
{\scriptsize\ttfamily NORMALIZED\_\allowbreak{}COMPLEMENT\_\allowbreak{}CRITERION} & {\scriptsize \#\# 6. Correct complement argument and compact action gap (lines 172-261)} & {\scriptsize proven} & {\scriptsize ---} \\
{\scriptsize\ttfamily ANGULAR\_\allowbreak{}REFERENCE\_\allowbreak{}BUDGET} & {\scriptsize \#\# 4. Antipodal geometry and a uniform angular reference calculation (lines 85-143)} & {\scriptsize proven} & {\scriptsize ---} \\
{\scriptsize\ttfamily ASSEMBLY\_\allowbreak{}CRITERION} & {\scriptsize \#\# 7. Assembly criterion and actual status (lines 262-286)} & {\scriptsize proven} & {\scriptsize ---} \\
{\scriptsize\ttfamily ACTUAL\_\allowbreak{}WEIGHTED\_\allowbreak{}JET\_\allowbreak{}AND\_\allowbreak{}LOWER\_\allowbreak{}BOUND} & {\scriptsize \#\# 6. Correct complement argument and compact action gap (lines 172-261)} & {\scriptsize proven} & {\scriptsize ---} \\
\end{longtable}

\subsection*{CITE:\allowbreak{}WILSON\_\allowbreak{}BACKGROUND}
\addcontentsline{toc}{subsection}{WILSON\_\allowbreak{}BACKGROUND}
\noindent\srcpath{docs/derivations/wilson-background-covariance-continuation.md}\\
{\footnotesize SHA-256 \texttt{05ee05dc5acf99fa}\ldots}\par\smallskip
\begin{longtable}{@{}p{0.30\linewidth} p{0.30\linewidth} l l@{}}
\toprule Statement & Locator & Status & Formal \\ \midrule
\endfirsthead
\toprule Statement & Locator & Status & Formal \\ \midrule \endhead
\bottomrule\endfoot
{\scriptsize\ttfamily BF1\_\allowbreak{}BF3} & {\scriptsize \#\# BF1. The energy-dependent relative-form statement (lines 48-102)} & {\scriptsize conditional} & {\scriptsize partial} \\
{\scriptsize\ttfamily BF4\_\allowbreak{}BF5} & {\scriptsize \#\# BF4. An actual Wilson small-field estimate, uniform in volume (lines 103-151)} & {\scriptsize proven} & {\scriptsize ---} \\
{\scriptsize\ttfamily BF6} & {\scriptsize \#\# BF6. Centering at an actual constrained minimizer (lines 152-188)} & {\scriptsize proven} & {\scriptsize ---} \\
{\scriptsize\ttfamily SC1\_\allowbreak{}SC4} & {\scriptsize \#\# 1. Setup and exact equations (lines 282-343)} & {\scriptsize proven} & {\scriptsize ---} \\
{\scriptsize\ttfamily SC5\_\allowbreak{}SC9} & {\scriptsize \#\# 2. A gauge-orbit gap that does not assume the physical gap (lines 344-415)} & {\scriptsize proven} & {\scriptsize ---} \\
{\scriptsize\ttfamily SC10\_\allowbreak{}SC15} & {\scriptsize \#\# 3. Uniform pointwise first and mixed derivatives at every coupling (lines 416-468)} & {\scriptsize proven} & {\scriptsize ---} \\
{\scriptsize\ttfamily SC16} & {\scriptsize \#\# 4. Direct finite-time form, without an implicit long-time interchange (lines 469-493)} & {\scriptsize proven} & {\scriptsize ---} \\
{\scriptsize\ttfamily SC17} & {\scriptsize \#\# 5. The neutral channel and the exact remaining spatial bound (lines 494-534)} & {\scriptsize conditional} & {\scriptsize ---} \\
{\scriptsize\ttfamily BC1} & {\scriptsize \#\# BC1. Why a fast inverse alone cannot finish the transport (lines 535-544)} & {\scriptsize proven} & {\scriptsize ---} \\
{\scriptsize\ttfamily BC2\_\allowbreak{}AUDIT} & {\scriptsize \#\# BC2. The additional multiscale document is evidence, not a completed input (lines 545-569)} & {\scriptsize proven} & {\scriptsize ---} \\
\end{longtable}

\subsection*{CITE:\allowbreak{}WILSON\_\allowbreak{}MARKED}
\addcontentsline{toc}{subsection}{WILSON\_\allowbreak{}MARKED}
\noindent\srcpath{docs/derivations/wilson-marked-transfer.md}\\
{\footnotesize SHA-256 \texttt{e7096d99ced64935}\ldots}\par\smallskip
\begin{longtable}{@{}p{0.30\linewidth} p{0.30\linewidth} l l@{}}
\toprule Statement & Locator & Status & Formal \\ \midrule
\endfirsthead
\toprule Statement & Locator & Status & Formal \\ \midrule \endhead
\bottomrule\endfoot
{\scriptsize\ttfamily WT1} & {\scriptsize \#\# WT-1. Exact physical-time blocks (lines 14-49)} & {\scriptsize proven} & {\scriptsize ---} \\
{\scriptsize\ttfamily WT2} & {\scriptsize \#\# WT-2. An all-orders activated-block bound (lines 50-88)} & {\scriptsize proven} & {\scriptsize ---} \\
{\scriptsize\ttfamily WT3} & {\scriptsize \#\# WT-3. Configurations and exact vacuum cancellation (lines 89-121)} & {\scriptsize proven} & {\scriptsize ---} \\
{\scriptsize\ttfamily WT4} & {\scriptsize \#\# WT-4. A uniform convergent vacuum polymer expansion (lines 122-171)} & {\scriptsize proven} & {\scriptsize ---} \\
{\scriptsize\ttfamily WT5} & {\scriptsize \#\# WT-5. Uniform control of free propagation and marked contour denominators (lines 172-208)} & {\scriptsize proven} & {\scriptsize ---} \\
{\scriptsize\ttfamily WT6} & {\scriptsize \#\# WT-6. The remaining complete-shell estimate and conditional transport (lines 209-248)} & {\scriptsize open} & {\scriptsize ---} \\
\end{longtable}

\subsection*{CITE:\allowbreak{}WILSON\_\allowbreak{}SC17}
\addcontentsline{toc}{subsection}{WILSON\_\allowbreak{}SC17}
\noindent\srcpath{docs/derivations/wilson-sc17-spatial-closure.md}\\
{\footnotesize SHA-256 \texttt{d02e906ba0d52d20}\ldots}\par\smallskip
\begin{longtable}{@{}p{0.30\linewidth} p{0.30\linewidth} l l@{}}
\toprule Statement & Locator & Status & Formal \\ \midrule
\endfirsthead
\toprule Statement & Locator & Status & Formal \\ \midrule \endhead
\bottomrule\endfoot
{\scriptsize\ttfamily REANCHORING} & {\scriptsize \#\# Endpoint reanchoring removes the neutral obstruction for norm bounds (lines 54-65)} & {\scriptsize proven} & {\scriptsize ---} \\
{\scriptsize\ttfamily CHARGED\_\allowbreak{}EVOLUTION} & {\scriptsize \#\# Why the charged contraction survives the time-dependent drift (lines 66-105)} & {\scriptsize proven} & {\scriptsize ---} \\
{\scriptsize\ttfamily H1\_\allowbreak{}H2} & {\scriptsize \#\# Actual weighted Volterra inequality (lines 106-131)} & {\scriptsize proven} & {\scriptsize ---} \\
{\scriptsize\ttfamily SMALL\_\allowbreak{}WINDOW} & {\scriptsize \#\# Explicit new regime (lines 132-166)} & {\scriptsize proven} & {\scriptsize ---} \\
{\scriptsize\ttfamily BARE\_\allowbreak{}MAJORANT\_\allowbreak{}WALL} & {\scriptsize \#\# Exact wall for this majorant (lines 167-189)} & {\scriptsize proven} & {\scriptsize ---} \\
{\scriptsize\ttfamily CURVATURE} & {\scriptsize \#\# Actual-vacuum curvature corollary (lines 190-212)} & {\scriptsize proven} & {\scriptsize ---} \\
{\scriptsize\ttfamily ANGLES} & {\scriptsize \#\# Sharper projection angles via covariance and the derived conditional gaps (lines 213-256)} & {\scriptsize proven} & {\scriptsize ---} \\
{\scriptsize\ttfamily OPTIMIZED\_\allowbreak{}INTERVAL} & {\scriptsize \#\# Maximized bare-Hessian bootstrap: a larger interval (lines 257-298)} & {\scriptsize proven} & {\scriptsize ---} \\
{\scriptsize\ttfamily RATIONAL\_\allowbreak{}WINDOW} & {\scriptsize \#\#\# Explicit rational larger-window corollary (lines 299-320)} & {\scriptsize proven} & {\scriptsize ---} \\
{\scriptsize\ttfamily GAUSSIAN\_\allowbreak{}PRECISION} & {\scriptsize \#\# 1. The actual flat quadratic ground precision (lines 323-358)} & {\scriptsize proven} & {\scriptsize ---} \\
{\scriptsize\ttfamily FRACTIONAL\_\allowbreak{}KERNEL} & {\scriptsize \#\# 2. Exact lattice fractional kernel and a sufficient weighted norm (lines 359-395)} & {\scriptsize proven} & {\scriptsize ---} \\
{\scriptsize\ttfamily GAUSSIAN\_\allowbreak{}ANGLE} & {\scriptsize \#\# 3. An exact scalar Gaussian angle test (lines 396-422)} & {\scriptsize proven} & {\scriptsize ---} \\
{\scriptsize\ttfamily FAST\_\allowbreak{}PRECISION\_\allowbreak{}CUSP} & {\scriptsize \#\# 4. Exact actual averaged-path fast conditional precision, ell=2 (lines 423-466)} & {\scriptsize proven} & {\scriptsize ---} \\
{\scriptsize\ttfamily R4A} & {\scriptsize \#\# 6. The existing dynamic covariance supplies the flat reference budget (lines 489-534)} & {\scriptsize proven} & {\scriptsize ---} \\
{\scriptsize\ttfamily R1\_\allowbreak{}R3} & {\scriptsize \#\# 1. Exact Euclidean equation and its projected cross terms (lines 537-573)} & {\scriptsize proven} & {\scriptsize ---} \\
{\scriptsize\ttfamily R4\_\allowbreak{}R9} & {\scriptsize \#\# 2. A rigorous conditional small-defect criterion (lines 574-648)} & {\scriptsize conditional} & {\scriptsize partial} \\
{\scriptsize\ttfamily R10\_\allowbreak{}R11} & {\scriptsize \#\# 3. An explicit true-angle consequence if the full Hessian defect is known (lines 649-697)} & {\scriptsize conditional} & {\scriptsize ---} \\
{\scriptsize\ttfamily R12\_\allowbreak{}R14} & {\scriptsize \#\# 4. The actual quantum pressure and cutoff defects are part of the input (lines 698-756)} & {\scriptsize conditional} & {\scriptsize ---} \\
{\scriptsize\ttfamily ACTUAL\_\allowbreak{}REFERENCE\_\allowbreak{}DEFECT\_\allowbreak{}BOUND} & {\scriptsize \#\# 4. The actual quantum pressure and cutoff defects are part of the input (lines 698-756), with the R5-R7 budget in section 2} & {\scriptsize open} & {\scriptsize ---} \\
{\scriptsize\ttfamily R15} & {\scriptsize \#\# 5. The averaged-path cusp is an analytic compression, not an eigenvalue crossing (lines 757-780)} & {\scriptsize proven} & {\scriptsize ---} \\
{\scriptsize\ttfamily PROPOSED\_\allowbreak{}QUANTUM\_\allowbreak{}PARAMETER\_\allowbreak{}WINDOW} & {\scriptsize \#\# 6. Formal repair of the spatial closure for lambda >= 1/73 via the quantum reference defect (lines 781-821)} & {\scriptsize conditional} & {\scriptsize partial} \\
\end{longtable}

\subsection*{CITE:\allowbreak{}WILSON\_\allowbreak{}SC17\_\allowbreak{}LIMIT}
\addcontentsline{toc}{subsection}{WILSON\_\allowbreak{}SC17\_\allowbreak{}LIMIT}
\noindent\srcpath{docs/derivations/wilson-sc17-thermodynamic-limit.md}\\
{\footnotesize SHA-256 \texttt{6d72876a35a86676}\ldots}\par\smallskip
\begin{longtable}{@{}p{0.30\linewidth} p{0.30\linewidth} l l@{}}
\toprule Statement & Locator & Status & Formal \\ \midrule
\endfirsthead
\toprule Statement & Locator & Status & Formal \\ \midrule \endhead
\bottomrule\endfoot
{\scriptsize\ttfamily T1} & {\scriptsize \#\# T1. Uniform setting including an actual cut interpolation (lines 28-95)} & {\scriptsize proven} & {\scriptsize ---} \\
{\scriptsize\ttfamily T2\_\allowbreak{}T3} & {\scriptsize \#\# T2. Weighted gradient contraction from the actual charged equations (lines 96-131)} & {\scriptsize proven} & {\scriptsize ---} \\
{\scriptsize\ttfamily T4\_\allowbreak{}T10} & {\scriptsize \#\# T3. The normalized actual vacuum response has no boundary-area loss (lines 132-195)} & {\scriptsize proven} & {\scriptsize ---} \\
{\scriptsize\ttfamily T11} & {\scriptsize \#\# T4. Actual exponential covariance localization (lines 196-210)} & {\scriptsize proven} & {\scriptsize ---} \\
{\scriptsize\ttfamily T12} & {\scriptsize \#\# T5. Fixed-lambda thermodynamic Cauchy convergence (lines 211-264)} & {\scriptsize proven} & {\scriptsize partial} \\
{\scriptsize\ttfamily T13} & {\scriptsize \#\# T7. Endpoint continuity also gives a thermodynamic limit at lambda\_\allowbreak{}c (lines 280-311)} & {\scriptsize proven} & {\scriptsize ---} \\
{\scriptsize\ttfamily SCORE\_\allowbreak{}LIMIT} & {\scriptsize \#\# Uniform score limits discharge the form construction's local input (lines 312-337)} & {\scriptsize proven} & {\scriptsize partial} \\
{\scriptsize\ttfamily IF3} & {\scriptsize \#\# 2. Exact infinite-volume integration by parts (lines 371-392)} & {\scriptsize proven} & {\scriptsize partial} \\
{\scriptsize\ttfamily IF4\_\allowbreak{}CLOSABLE} & {\scriptsize \#\# 3. Closability and the canonical closed form (lines 393-436)} & {\scriptsize proven} & {\scriptsize partial} \\
{\scriptsize\ttfamily IF4\_\allowbreak{}IF7} & {\scriptsize \#\# 3. Closability and the canonical closed form (lines 393-436)} & {\scriptsize proven} & {\scriptsize partial} \\
{\scriptsize\ttfamily PHYSICAL\_\allowbreak{}RESTRICTION} & {\scriptsize \#\# 4. Gauge invariance and the physical closed subspace (lines 437-460)} & {\scriptsize proven} & {\scriptsize ---} \\
{\scriptsize\ttfamily IF8\_\allowbreak{}LIMIT} & {\scriptsize \#\# 5. Passing the full and physical Poincare floor (lines 461-479)} & {\scriptsize proven} & {\scriptsize partial} \\
{\scriptsize\ttfamily IF8\_\allowbreak{}IF9} & {\scriptsize \#\# 5. Passing the full and physical Poincare floor (lines 461-479)} & {\scriptsize proven} & {\scriptsize partial} \\
{\scriptsize\ttfamily IF10\_\allowbreak{}IF13} & {\scriptsize \#\# 6. Actual nonzero local observable overlap with a bounded energy band (lines 480-528)} & {\scriptsize proven} & {\scriptsize partial} \\
\end{longtable}

\subsection*{CITE:\allowbreak{}WILSON\_\allowbreak{}SC17\_\allowbreak{}TIME}
\addcontentsline{toc}{subsection}{WILSON\_\allowbreak{}SC17\_\allowbreak{}TIME}
\noindent\srcpath{docs/derivations/wilson-sc17-physical-time-limit.md}\\
{\footnotesize SHA-256 \texttt{f842a6b9f826a997}\ldots}\par\smallskip
\begin{longtable}{@{}p{0.30\linewidth} p{0.30\linewidth} l l@{}}
\toprule Statement & Locator & Status & Formal \\ \midrule
\endfirsthead
\toprule Statement & Locator & Status & Formal \\ \midrule \endhead
\bottomrule\endfoot
{\scriptsize\ttfamily P4\_\allowbreak{}P6} & {\scriptsize \#\# P2. A limiting influence kernel and weighted Lipschitz estimate (lines 75-133)} & {\scriptsize proven} & {\scriptsize ---} \\
{\scriptsize\ttfamily P7\_\allowbreak{}P13} & {\scriptsize \#\# P3. A global extension with exact sphere invariance (lines 134-220)} & {\scriptsize proven} & {\scriptsize ---} \\
{\scriptsize\ttfamily P14} & {\scriptsize \#\# P4. Uniform Hilbert convergence of the actual finite drifts (lines 221-258)} & {\scriptsize proven} & {\scriptsize ---} \\
{\scriptsize\ttfamily P15} & {\scriptsize \#\# P5. Strong finite-time process comparison (lines 259-293)} & {\scriptsize proven} & {\scriptsize ---} \\
{\scriptsize\ttfamily P16} & {\scriptsize \#\# P6. Stationary path laws and actual multitime vacuum correlations (lines 294-332)} & {\scriptsize proven} & {\scriptsize ---} \\
{\scriptsize\ttfamily P17\_\allowbreak{}P18} & {\scriptsize \#\# P7. Identification with the already closed gradient form (lines 333-440)} & {\scriptsize proven} & {\scriptsize ---} \\
{\scriptsize\ttfamily PHYSICAL\_\allowbreak{}TIME\_\allowbreak{}GAP} & {\scriptsize \#\# P8. Physical restriction and the fixed-spacing time limit (lines 441-470)} & {\scriptsize proven} & {\scriptsize ---} \\
{\scriptsize\ttfamily P19\_\allowbreak{}P20} & {\scriptsize \#\# P9. Endpoint extension by a summable nonexponential weight (lines 471-525)} & {\scriptsize proven} & {\scriptsize ---} \\
\end{longtable}

\subsection*{CITE:\allowbreak{}WILSON\_\allowbreak{}SELECTED}
\addcontentsline{toc}{subsection}{WILSON\_\allowbreak{}SELECTED}
\noindent\srcpath{docs/derivations/wilson-selected-inverse-wall.md}\\
{\footnotesize SHA-256 \texttt{c0473aca76790ef7}\ldots}\par\smallskip
\begin{longtable}{@{}p{0.30\linewidth} p{0.30\linewidth} l l@{}}
\toprule Statement & Locator & Status & Formal \\ \midrule
\endfirsthead
\toprule Statement & Locator & Status & Formal \\ \midrule \endhead
\bottomrule\endfoot
{\scriptsize\ttfamily W1\_\allowbreak{}W3\_\allowbreak{}OBSTRUCTION} & {\scriptsize \#\# 2. First attempted step: global relative Gaussian fast-form control (lines 33-100)} & {\scriptsize proven} & {\scriptsize ---} \\
{\scriptsize\ttfamily W4\_\allowbreak{}W5} & {\scriptsize \#\# 3. Repair: only compare inverses on the graph force (lines 101-157)} & {\scriptsize conditional} & {\scriptsize partial} \\
{\scriptsize\ttfamily SIGNED\_\allowbreak{}PAIRING\_\allowbreak{}OBSTRUCTIONS} & {\scriptsize \#\# 4. Attempt to establish the repaired estimate for Wilson (lines 158-220)} & {\scriptsize proven} & {\scriptsize ---} \\
{\scriptsize\ttfamily W6} & {\scriptsize \#\# 5. The exact point at which this argument cannot proceed (lines 221-264)} & {\scriptsize open} & {\scriptsize partial} \\
\end{longtable}

\subsection*{CITE:\allowbreak{}WILSON\_\allowbreak{}SHELL}
\addcontentsline{toc}{subsection}{WILSON\_\allowbreak{}SHELL}
\noindent\srcpath{docs/derivations/wilson-marked-shell-transport.md}\\
{\footnotesize SHA-256 \texttt{a887ed9de154cec0}\ldots}\par\smallskip
\begin{longtable}{@{}p{0.30\linewidth} p{0.30\linewidth} l l@{}}
\toprule Statement & Locator & Status & Formal \\ \midrule
\endfirsthead
\toprule Statement & Locator & Status & Formal \\ \midrule \endhead
\bottomrule\endfoot
{\scriptsize\ttfamily S2\_\allowbreak{}S7} & {\scriptsize \#\# 2. The expansion with a fixed observable (lines 76-169)} & {\scriptsize proven} & {\scriptsize ---} \\
{\scriptsize\ttfamily S8\_\allowbreak{}S9} & {\scriptsize \#\# 3. Complex anchored operator bound (lines 170-218)} & {\scriptsize proven} & {\scriptsize ---} \\
{\scriptsize\ttfamily S10\_\allowbreak{}S15} & {\scriptsize \#\# 4. Actual complete shell and a holomorphic literal-source frame (lines 219-311)} & {\scriptsize proven} & {\scriptsize ---} \\
{\scriptsize\ttfamily S16\_\allowbreak{}S17} & {\scriptsize \#\# 5. From finite-order locality to a uniform weighted Taylor bound (lines 312-390)} & {\scriptsize proven} & {\scriptsize ---} \\
{\scriptsize\ttfamily S18\_\allowbreak{}S20} & {\scriptsize \#\# 6. Sum the matching and transport the carrier (lines 391-461)} & {\scriptsize conditional} & {\scriptsize ---} \\
{\scriptsize\ttfamily S21\_\allowbreak{}S22} & {\scriptsize \#\# 7. Observable overlap, with the quantifiers retained (lines 462-493)} & {\scriptsize proven} & {\scriptsize ---} \\
\end{longtable}

\subsection*{CITE:\allowbreak{}WILSON\_\allowbreak{}SPATIAL}
\addcontentsline{toc}{subsection}{WILSON\_\allowbreak{}SPATIAL}
\noindent\srcpath{docs/derivations/wilson-spatial-schur-excess.md}\\
{\footnotesize SHA-256 \texttt{5ddb7f812a9e6f53}\ldots}\par\smallskip
\begin{longtable}{@{}p{0.30\linewidth} p{0.30\linewidth} l l@{}}
\toprule Statement & Locator & Status & Formal \\ \midrule
\endfirsthead
\toprule Statement & Locator & Status & Formal \\ \midrule \endhead
\bottomrule\endfoot
{\scriptsize\ttfamily SP1\_\allowbreak{}SP6} & {\scriptsize \#\# 2. Exact excess about the full reference graph (lines 42-114)} & {\scriptsize conditional} & {\scriptsize partial} \\
{\scriptsize\ttfamily SP7\_\allowbreak{}SP10} & {\scriptsize \#\# 3. The complete second-order coefficient and a uniform remainder (lines 115-182)} & {\scriptsize conditional} & {\scriptsize partial} \\
{\scriptsize\ttfamily SP11\_\allowbreak{}SP16} & {\scriptsize \#\# 4. What W\_\allowbreak{}1 and W\_\allowbreak{}2 mean for the actual Wilson operator (lines 183-306)} & {\scriptsize proven} & {\scriptsize ---} \\
{\scriptsize\ttfamily SP17\_\allowbreak{}SP19} & {\scriptsize \#\# 5. Second-order literal-source normalization and harmonic control (lines 307-411)} & {\scriptsize proven} & {\scriptsize ---} \\
{\scriptsize\ttfamily SP20\_\allowbreak{}SP24} & {\scriptsize \#\# 6. Iterating energies and observable amplitudes in the physical clock (lines 412-477)} & {\scriptsize conditional} & {\scriptsize ---} \\
{\scriptsize\ttfamily SP25} & {\scriptsize \#\# 7. Why the complete second-order subtraction matters (lines 478-519)} & {\scriptsize proven} & {\scriptsize ---} \\
\end{longtable}

\subsection*{CITE:\allowbreak{}WILSON\_\allowbreak{}TRUE\_\allowbreak{}BLOCK}
\addcontentsline{toc}{subsection}{WILSON\_\allowbreak{}TRUE\_\allowbreak{}BLOCK}
\noindent\srcpath{docs/derivations/wilson-true-vacuum-block-estimates.md}\\
{\footnotesize SHA-256 \texttt{23aebb949cd2b6da}\ldots}\par\smallskip
\begin{longtable}{@{}p{0.30\linewidth} p{0.30\linewidth} l l@{}}
\toprule Statement & Locator & Status & Formal \\ \midrule
\endfirsthead
\toprule Statement & Locator & Status & Formal \\ \midrule \endhead
\bottomrule\endfoot
{\scriptsize\ttfamily BA1} & {\scriptsize \#\# BA1. A specified block partition (lines 10-28)} & {\scriptsize proven} & {\scriptsize ---} \\
{\scriptsize\ttfamily BA2\_\allowbreak{}BA5} & {\scriptsize \#\# BA2. Optimizing the true conditional gap (lines 29-94)} & {\scriptsize proven} & {\scriptsize partial} \\
{\scriptsize\ttfamily BA6\_\allowbreak{}BA8} & {\scriptsize \#\# BA3. Reduce the actual projection angle to a mixed ground-amplitude ratio (lines 95-141)} & {\scriptsize proven} & {\scriptsize partial} \\
{\scriptsize\ttfamily BA9\_\allowbreak{}BA10} & {\scriptsize \#\# BA4. What a frozen-block spectral computation would additionally require (lines 142-168)} & {\scriptsize conditional} & {\scriptsize ---} \\
{\scriptsize\ttfamily BA11\_\allowbreak{}BA12} & {\scriptsize \#\# BA5. Exact finite-time derivation of the connected quantity (lines 169-201)} & {\scriptsize proven} & {\scriptsize ---} \\
{\scriptsize\ttfamily BA13\_\allowbreak{}BA16} & {\scriptsize \#\# BA6. The first connected coefficients for the actual coupled lattice (lines 202-269)} & {\scriptsize proven} & {\scriptsize ---} \\
{\scriptsize\ttfamily BA20\_\allowbreak{}BA26} & {\scriptsize \#\# BA7. Actual angle closure in an explicit strong-coupling window (lines 270-405)} & {\scriptsize proven} & {\scriptsize ---} \\
\end{longtable}

\subsection*{CITE:\allowbreak{}WILSON\_\allowbreak{}TRUE\_\allowbreak{}BLOCK\_\allowbreak{}CHECK}
\addcontentsline{toc}{subsection}{WILSON\_\allowbreak{}TRUE\_\allowbreak{}BLOCK\_\allowbreak{}CHECK}
\noindent\srcpath{docs/derivations/wilson-true-blocks-independent-check.md}\\
{\footnotesize SHA-256 \texttt{6984f428f801c70e}\ldots}\par\smallskip
\begin{longtable}{@{}p{0.30\linewidth} p{0.30\linewidth} l l@{}}
\toprule Statement & Locator & Status & Formal \\ \midrule
\endfirsthead
\toprule Statement & Locator & Status & Formal \\ \midrule \endhead
\bottomrule\endfoot
{\scriptsize\ttfamily ENVELOPE\_\allowbreak{}AND\_\allowbreak{}CROSS\_\allowbreak{}RATIO} & {\scriptsize \#\# CB1. What was re-derived independently (lines 17-56)} & {\scriptsize proven} & {\scriptsize ---} \\
{\scriptsize\ttfamily CB1} & {\scriptsize \#\# CB2. The single-link partition is stronger, and closes VA12 directly (lines 57-75)} & {\scriptsize proven} & {\scriptsize ---} \\
{\scriptsize\ttfamily FREE\_\allowbreak{}BUDGET\_\allowbreak{}LOOSENESS} & {\scriptsize \#\# CB3. The angle budget is not graded in lambda (lines 76-87)} & {\scriptsize proven} & {\scriptsize ---} \\
{\scriptsize\ttfamily CB2} & {\scriptsize \#\# CB4. The BA26 wall is more robust than its own derivation (lines 88-109)} & {\scriptsize proven} & {\scriptsize ---} \\
{\scriptsize\ttfamily NONCHAINING} & {\scriptsize \#\# CB5. The two boxed constants must not be chained (lines 110-121)} & {\scriptsize proven} & {\scriptsize ---} \\
\end{longtable}

\subsection*{CITE:\allowbreak{}WILSON\_\allowbreak{}VACUUM\_\allowbreak{}ASSEMBLY}
\addcontentsline{toc}{subsection}{WILSON\_\allowbreak{}VACUUM\_\allowbreak{}ASSEMBLY}
\noindent\srcpath{docs/derivations/wilson-vacuum-aligned-assembly.md}\\
{\footnotesize SHA-256 \texttt{6c21bda892b88872}\ldots}\par\smallskip
\begin{longtable}{@{}p{0.30\linewidth} p{0.30\linewidth} l l@{}}
\toprule Statement & Locator & Status & Formal \\ \midrule
\endfirsthead
\toprule Statement & Locator & Status & Formal \\ \midrule \endhead
\bottomrule\endfoot
{\scriptsize\ttfamily VA1\_\allowbreak{}VA4} & {\scriptsize \#\# VA1. A local finite-energy bound for the actual quantum vacuum (lines 17-74)} & {\scriptsize proven} & {\scriptsize partial} \\
{\scriptsize\ttfamily VA5\_\allowbreak{}VA7} & {\scriptsize \#\# VA2. The WR25 remainder is extensive on actual connected lattices (lines 75-130)} & {\scriptsize proven} & {\scriptsize partial} \\
{\scriptsize\ttfamily VA8} & {\scriptsize \#\# VA3. WR26 is false on the physical vacuum complement at large volume (lines 131-164)} & {\scriptsize proven} & {\scriptsize ---} \\
{\scriptsize\ttfamily VA9\_\allowbreak{}VA11} & {\scriptsize \#\# VA4. Exact cancellation before comparison (lines 165-210)} & {\scriptsize proven} & {\scriptsize partial} \\
{\scriptsize\ttfamily VA12\_\allowbreak{}VA13} & {\scriptsize \#\# VA5. Projection geometry, with the Gauss law retained (lines 211-251)} & {\scriptsize conditional} & {\scriptsize partial} \\
{\scriptsize\ttfamily VA14\_\allowbreak{}VA19} & {\scriptsize \#\# VA6. The next estimate and the attempts to establish it (lines 252-415)} & {\scriptsize conditional} & {\scriptsize partial} \\
\end{longtable}

\subsection*{CITE:\allowbreak{}WILSON\_\allowbreak{}WEIGHTED\_\allowbreak{}REPAIR}
\addcontentsline{toc}{subsection}{WILSON\_\allowbreak{}WEIGHTED\_\allowbreak{}REPAIR}
\noindent\srcpath{docs/derivations/wilson-weighted-repair-and-rotor-gap.md}\\
{\footnotesize SHA-256 \texttt{dac4e8fd2d13f3e7}\ldots}\par\smallskip
\begin{longtable}{@{}p{0.30\linewidth} p{0.30\linewidth} l l@{}}
\toprule Statement & Locator & Status & Formal \\ \midrule
\endfirsthead
\toprule Statement & Locator & Status & Formal \\ \midrule \endhead
\bottomrule\endfoot
{\scriptsize\ttfamily WR1\_\allowbreak{}WR3} & {\scriptsize \#\# WR1. Absorb negative curvature in the derivative energy (lines 15-66)} & {\scriptsize proven} & {\scriptsize ---} \\
{\scriptsize\ttfamily WR4\_\allowbreak{}WR6} & {\scriptsize \#\# WR2. A smooth non-Gaussian trial for the physical rotor (lines 67-99)} & {\scriptsize proven} & {\scriptsize ---} \\
{\scriptsize\ttfamily WR7\_\allowbreak{}WR12} & {\scriptsize \#\# WR3. A physical gap on the complete rotor space, at every coupling (lines 100-161)} & {\scriptsize proven} & {\scriptsize ---} \\
{\scriptsize\ttfamily ROTOR\_\allowbreak{}VOLUME} & {\scriptsize \#\# WR4. Volume passage that is already justified (lines 162-175)} & {\scriptsize proven} & {\scriptsize ---} \\
{\scriptsize\ttfamily WR13\_\allowbreak{}WR16} & {\scriptsize \#\# WR5. Actual shared-link residual, and a positive interacting repair (lines 176-228)} & {\scriptsize proven} & {\scriptsize ---} \\
{\scriptsize\ttfamily WR17\_\allowbreak{}WR21} & {\scriptsize \#\# WR6. The precise remaining derivation (lines 229-308)} & {\scriptsize conditional} & {\scriptsize ---} \\
{\scriptsize\ttfamily WR22\_\allowbreak{}WR24} & {\scriptsize \#\# WR7. Retain the full boundary: an actual-link fast barrier (lines 309-356)} & {\scriptsize proven} & {\scriptsize ---} \\
{\scriptsize\ttfamily WR25\_\allowbreak{}WR28} & {\scriptsize \#\# WR8. Where the repaired local bounds still require a new proof (lines 357-443)} & {\scriptsize proven} & {\scriptsize ---} \\
\end{longtable}

\subsection*{CITE:\allowbreak{}YM\_\allowbreak{}BALABAN\_\allowbreak{}MULTISCALE}
\addcontentsline{toc}{subsection}{YM\_\allowbreak{}BALABAN\_\allowbreak{}MULTISCALE}
\noindent\srcpath{docs/derivations/yangmills-continuum-balaban-multiscale-proof.md}\\
{\footnotesize SHA-256 \texttt{55296902489318e2}\ldots}\par\smallskip
\begin{longtable}{@{}p{0.30\linewidth} p{0.30\linewidth} l l@{}}
\toprule Statement & Locator & Status & Formal \\ \midrule
\endfirsthead
\toprule Statement & Locator & Status & Formal \\ \midrule \endhead
\bottomrule\endfoot
{\scriptsize\ttfamily LEMMA\_\allowbreak{}3\_\allowbreak{}2\_\allowbreak{}ADJOINT} & {\scriptsize \#\#\# 3.1 Geometric Covariant Averaging Operator \$Q\_\allowbreak{}k\$ (lines 84-110)} & {\scriptsize proven} & {\scriptsize ---} \\
{\scriptsize\ttfamily PROPOSITION\_\allowbreak{}3\_\allowbreak{}4} & {\scriptsize \#\#\# 3.2 The Background Field Variational Problem (lines 111-139)} & {\scriptsize disputed} & {\scriptsize ---} \\
{\scriptsize\ttfamily THEOREM\_\allowbreak{}4\_\allowbreak{}2} & {\scriptsize \#\#\# 4.2 Super-Exponential Suppression of Large Fields (lines 157-180)} & {\scriptsize conditional} & {\scriptsize ---} \\
{\scriptsize\ttfamily GAUGE\_\allowbreak{}HESSIAN} & {\scriptsize \#\#\# 5.1 Fluctuation Gauge Fixing (lines 183-197)} & {\scriptsize disputed} & {\scriptsize ---} \\
{\scriptsize\ttfamily GHOST\_\allowbreak{}5\_\allowbreak{}8} & {\scriptsize \#\#\# 5.2 Faddeev--Popov Ghost Determinant (lines 198-208)} & {\scriptsize conditional} & {\scriptsize ---} \\
{\scriptsize\ttfamily THEOREM\_\allowbreak{}5\_\allowbreak{}1} & {\scriptsize \#\#\# 5.3 Combes--Thomas Resolvent Localization for Continuous \$\textbackslash{}mathcal\{H\}\_\allowbreak{}\{B\_\allowbreak{}k\}\$ (lines 209-248)} & {\scriptsize disputed} & {\scriptsize ---} \\
{\scriptsize\ttfamily BANACH\_\allowbreak{}ALGEBRA} & {\scriptsize \#\#\# 6.2 The Polymer Banach Algebra \$\textbackslash{}mathcal\{E\}\_\allowbreak{}k\$ (lines 260-268)} & {\scriptsize proven} & {\scriptsize ---} \\
{\scriptsize\ttfamily THEOREM\_\allowbreak{}6\_\allowbreak{}2} & {\scriptsize \#\#\# 6.3 The Mayer Tree-Graph Expansion \& Cluster Contraction (lines 269-295)} & {\scriptsize conditional} & {\scriptsize ---} \\
{\scriptsize\ttfamily THEOREM\_\allowbreak{}6\_\allowbreak{}3} & {\scriptsize \#\#\# 6.4 Asymptotic Freedom and Infinite Scale Stability (lines 296-314)} & {\scriptsize conditional} & {\scriptsize ---} \\
{\scriptsize\ttfamily THEOREM\_\allowbreak{}7\_\allowbreak{}1} & {\scriptsize \#\#\# 7.1 Cauchy Sequence of Effective Actions (lines 317-340)} & {\scriptsize open} & {\scriptsize ---} \\
{\scriptsize\ttfamily THEOREM\_\allowbreak{}7\_\allowbreak{}2} & {\scriptsize \#\#\# 7.2 Construction of the Limiting Continuum Measure \$d\textbackslash{}mu\_\allowbreak{}\{\textbackslash{}text\{YM\}\}\$ (lines 341-363)} & {\scriptsize conditional} & {\scriptsize ---} \\
{\scriptsize\ttfamily OS0} & {\scriptsize \#\#\# 8.1 OS0: Analyticity and Temperedness (lines 368-374)} & {\scriptsize conditional} & {\scriptsize ---} \\
{\scriptsize\ttfamily OS1} & {\scriptsize \#\#\# 8.2 OS1: Euclidean Covariance (lines 375-382)} & {\scriptsize open} & {\scriptsize ---} \\
{\scriptsize\ttfamily THEOREM\_\allowbreak{}8\_\allowbreak{}1} & {\scriptsize \#\#\# 8.3 OS2: Reflection Positivity (lines 383-408)} & {\scriptsize conditional} & {\scriptsize ---} \\
{\scriptsize\ttfamily OS3} & {\scriptsize \#\#\# 8.4 OS3: Permutation Symmetry (lines 409-411)} & {\scriptsize conditional} & {\scriptsize ---} \\
{\scriptsize\ttfamily THEOREM\_\allowbreak{}8\_\allowbreak{}2} & {\scriptsize \#\#\# 8.5 OS4: Exponential Clustering and Mass Gap (lines 412-436)} & {\scriptsize open} & {\scriptsize ---} \\
{\scriptsize\ttfamily OS\_\allowbreak{}RECONSTRUCTION} & {\scriptsize \#\#\# 9.1 The Physical Hilbert Space and Hamiltonian (lines 439-453)} & {\scriptsize conditional} & {\scriptsize ---} \\
{\scriptsize\ttfamily THEOREM\_\allowbreak{}9\_\allowbreak{}1} & {\scriptsize \#\#\# 9.2 The Physical Spectral Mass Gap (lines 454-478)} & {\scriptsize conditional} & {\scriptsize ---} \\
{\scriptsize\ttfamily SCATTERING} & {\scriptsize \#\#\# 9.3 Relativistic Wightman Fields and Non-Trivial Scattering (\$S \textbackslash{}ne I\$) (lines 479-501)} & {\scriptsize conditional} & {\scriptsize ---} \\
\end{longtable}

\subsection*{CITE:\allowbreak{}YM\_\allowbreak{}FLAT\_\allowbreak{}DIRECTIONS}
\addcontentsline{toc}{subsection}{YM\_\allowbreak{}FLAT\_\allowbreak{}DIRECTIONS}
\noindent\srcpath{docs/derivations/yangmills-simon-flat-directions.md}\\
{\footnotesize SHA-256 \texttt{850d8bf0ea4b6192}\ldots}\par\smallskip
\begin{longtable}{@{}p{0.30\linewidth} p{0.30\linewidth} l l@{}}
\toprule Statement & Locator & Status & Formal \\ \midrule
\endfirsthead
\toprule Statement & Locator & Status & Formal \\ \midrule \endhead
\bottomrule\endfoot
{\scriptsize\ttfamily OSCILLATOR} & {\scriptsize \#\# 1. The scalar oscillator, including its constant (lines 102-138)} & {\scriptsize proven} & {\scriptsize ---} \\
{\scriptsize\ttfamily TRANSVERSE} & {\scriptsize \#\# 2. SU(2) has two transverse oscillator modes per slice (lines 139-167)} & {\scriptsize proven} & {\scriptsize ---} \\
{\scriptsize\ttfamily SLICE\_\allowbreak{}AGGREGATION} & {\scriptsize \#\# 3. Allocate the pair potential without losing longitudinal control (lines 168-212)} & {\scriptsize proven} & {\scriptsize ---} \\
{\scriptsize\ttfamily COMPACT\_\allowbreak{}RESOLVENT} & {\scriptsize \#\# 4. Compact resolvent follows from the retained gradient and radial tail (lines 213-233)} & {\scriptsize proven} & {\scriptsize ---} \\
{\scriptsize\ttfamily GROUND\_\allowbreak{}FLOOR} & {\scriptsize \#\# 5. An explicit ground-energy lower bound (lines 234-275)} & {\scriptsize proven} & {\scriptsize ---} \\
{\scriptsize\ttfamily JACOBIAN\_\allowbreak{}VALLEY} & {\scriptsize \#\# 6. What Simon's Jacobian condition measures (lines 276-324)} & {\scriptsize proven} & {\scriptsize ---} \\
{\scriptsize\ttfamily SCALING} & {\scriptsize \#\# 7. Energy scale, vacuum subtraction and the research gaps (lines 325-368)} & {\scriptsize proven} & {\scriptsize ---} \\
\end{longtable}

\subsection*{CITE:\allowbreak{}YM\_\allowbreak{}GPU\_\allowbreak{}AUDIT}
\addcontentsline{toc}{subsection}{YM\_\allowbreak{}GPU\_\allowbreak{}AUDIT}
\noindent\srcpath{docs/derivations/yangmills-gpu-resolution-audit.md}\\
{\footnotesize SHA-256 \texttt{45a72f1f90d9da8f}\ldots}\par\smallskip
\begin{longtable}{@{}p{0.30\linewidth} p{0.30\linewidth} l l@{}}
\toprule Statement & Locator & Status & Formal \\ \midrule
\endfirsthead
\toprule Statement & Locator & Status & Formal \\ \midrule \endhead
\bottomrule\endfoot
{\scriptsize\ttfamily GA1\_\allowbreak{}GA4} & {\scriptsize \#\# GA1. Phase 1 is a massive scalar model with an exact uniform theorem (lines 24-76)} & {\scriptsize proven} & {\scriptsize ---} \\
{\scriptsize\ttfamily GA5\_\allowbreak{}GA7} & {\scriptsize \#\# GA2. Phase 2's alleged domino is an exact tensor sum (lines 77-129)} & {\scriptsize proven} & {\scriptsize ---} \\
{\scriptsize\ttfamily RADIAL\_\allowbreak{}DISCRETIZATION} & {\scriptsize \#\# GA3. Phase 3 compares two radial discretizations, not a transfer matrix (lines 130-161)} & {\scriptsize proven} & {\scriptsize ---} \\
{\scriptsize\ttfamily GA8\_\allowbreak{}GA9} & {\scriptsize \#\# GA4. The additional Phase 4 has a valid conditional scalar limit (lines 162-209)} & {\scriptsize conditional} & {\scriptsize ---} \\
{\scriptsize\ttfamily GA10} & {\scriptsize \#\# GA5. Phase 5 keeps an admissible cube basis but changes its magnetic operator (lines 210-249)} & {\scriptsize proven} & {\scriptsize partial} \\
{\scriptsize\ttfamily GA11} & {\scriptsize \#\# GA6. What the supplied Lean additions actually formalize (lines 250-286)} & {\scriptsize proven} & {\scriptsize ---} \\
{\scriptsize\ttfamily GA12} & {\scriptsize \#\# GA7. The incoming smearing script measures spatial power, not spectral overlap (lines 287-327)} & {\scriptsize proven} & {\scriptsize ---} \\
{\scriptsize\ttfamily GA13} & {\scriptsize \#\# GA8. The later G17 script inserts both its cluster coefficients and its mass (lines 328-363)} & {\scriptsize proven} & {\scriptsize ---} \\
{\scriptsize\ttfamily GA14\_\allowbreak{}GA19} & {\scriptsize \#\# GA9. Phase 8 supplies an assumed block kernel, not conditional projections (lines 364-500)} & {\scriptsize proven} & {\scriptsize ---} \\
{\scriptsize\ttfamily GA20\_\allowbreak{}GA21} & {\scriptsize \#\# GA10. Phase 9 verifies a Gram identity, an inserted exponential and a Haar moment (lines 501-563)} & {\scriptsize proven} & {\scriptsize partial} \\
\end{longtable}

\subsection*{CITE:\allowbreak{}YM\_\allowbreak{}GROUND\_\allowbreak{}STATE}
\addcontentsline{toc}{subsection}{YM\_\allowbreak{}GROUND\_\allowbreak{}STATE}
\noindent\srcpath{docs/derivations/yangmills-weighted-curvature.md}\\
{\footnotesize SHA-256 \texttt{75c853725c68db30}\ldots}\par\smallskip
\begin{longtable}{@{}p{0.30\linewidth} p{0.30\linewidth} l l@{}}
\toprule Statement & Locator & Status & Formal \\ \midrule
\endfirsthead
\toprule Statement & Locator & Status & Formal \\ \midrule \endhead
\bottomrule\endfoot
{\scriptsize\ttfamily GST1} & {\scriptsize \#\# GST-1: exact conjugation with a trial-state residual (lines 32-70)} & {\scriptsize proven} & {\scriptsize partial} \\
{\scriptsize\ttfamily GST2} & {\scriptsize \#\# GST-2: weighted energy and the needed domain statement (lines 71-91)} & {\scriptsize proven} & {\scriptsize partial} \\
{\scriptsize\ttfamily GST3} & {\scriptsize \#\# GST-3: curvature normalization and a spectral argument without a discrete excited state (lines 92-130)} & {\scriptsize conditional} & {\scriptsize partial} \\
{\scriptsize\ttfamily GST4} & {\scriptsize \#\# GST-4: the finite matrix identity and its exact hypotheses (lines 131-146)} & {\scriptsize proven} & {\scriptsize full} \\
{\scriptsize\ttfamily GST5} & {\scriptsize \#\# GST-5: oscillator calibration (lines 147-166)} & {\scriptsize proven} & {\scriptsize ---} \\
{\scriptsize\ttfamily GST6} & {\scriptsize \#\# GST-6: the residual-mean certificate (lines 167-190)} & {\scriptsize conditional} & {\scriptsize ---} \\
{\scriptsize\ttfamily GST7} & {\scriptsize \#\# GST-7: a complete algebraic certificate for the pure quartic oscillator (lines 191-244)} & {\scriptsize proven} & {\scriptsize ---} \\
{\scriptsize\ttfamily GST8} & {\scriptsize \#\# GST-8: falsifiers and the Yang--Mills boundary (lines 245-266)} & {\scriptsize proven} & {\scriptsize ---} \\
{\scriptsize\ttfamily GST9} & {\scriptsize \#\# GST-9: an exact volume-accumulation falsifier (lines 267-300)} & {\scriptsize proven} & {\scriptsize ---} \\
\end{longtable}

\subsection*{CITE:\allowbreak{}YM\_\allowbreak{}RECONSTRUCTION}
\addcontentsline{toc}{subsection}{YM\_\allowbreak{}RECONSTRUCTION}
\noindent\srcpath{docs/derivations/yangmills-reconstruction.md}\\
{\footnotesize SHA-256 \texttt{a180dddf14179838}\ldots}\par\smallskip
\begin{longtable}{@{}p{0.30\linewidth} p{0.30\linewidth} l l@{}}
\toprule Statement & Locator & Status & Formal \\ \midrule
\endfirsthead
\toprule Statement & Locator & Status & Formal \\ \midrule \endhead
\bottomrule\endfoot
{\scriptsize\ttfamily R1} & {\scriptsize \#\# 1. Source statements and the operator to be bounded (lines 5-27)} & {\scriptsize proven} & {\scriptsize partial} \\
{\scriptsize\ttfamily R2} & {\scriptsize \#\# 2. A finite reflection-positive model one can inspect exactly (lines 28-38)} & {\scriptsize proven} & {\scriptsize full} \\
{\scriptsize\ttfamily R3} & {\scriptsize \#\# 3. Observable completeness is a testable interface (lines 39-51)} & {\scriptsize proven} & {\scriptsize full} \\
{\scriptsize\ttfamily R4\_\allowbreak{}R5} & {\scriptsize \#\# 4. What a localization plateau actually proves (lines 52-89)} & {\scriptsize proven} & {\scriptsize partial} \\
{\scriptsize\ttfamily R6\_\allowbreak{}R7} & {\scriptsize \#\# 5. A precise sufficient repair, with its remaining hypothesis exposed (lines 90-111)} & {\scriptsize conditional} & {\scriptsize partial} \\
\end{longtable}



\section{Statements that do not close}
\label{app:open}

\DerivOpen{} of the \DerivStatements{} registered statements are
recorded as open, conditional, or disputed. They are extracted here
with their hypotheses and the corpus's own sentence on what remains,
because this is the shortest complete list of places where an agent can
do mathematics that the program has already established is missing.

Read alongside \S\ref{sec:obligations}, which curates the subset the
program judges decisive, and \S\ref{sec:attack}, which is about how to
start on one.

% Generated by paper/master/generate_sources.py.  Do not edit.
\noindent 63 of the registered derivation statements are not recorded as proven. Each is listed with the document it lives in, its locator, and the corpus's own sentence on what remains. These are the statements an agent can act on directly.\par\medskip
\begingroup\sloppy
\noindent\textbf{\texttt{\small DERIV:\allowbreak{}EARLIER\_\allowbreak{}PROOF\_\allowbreak{}RECOVERY:\allowbreak{}DIPOLAR\_\allowbreak{}REMAINDER}} \hfill {\footnotesize open}\par
{\small Establish the archived O(|x|\textasciicircum{}-4) remainder of the dipolar lattice kernel with the required Fourier-cutoff and derivative bounds.}\par
{\footnotesize \srcpath{docs/derivations/earlier-proof-recovery.md} --- \#\# E18. Explicit remaining application obligations}\par
{\footnotesize \textit{Hypotheses.} The exact source assumptions and required estimates described in E18.}\par
{\footnotesize \textit{Remaining.} Explicit pending application or extension; neither the exact finite controls nor the narrower recovered theorem discharge it.}\par
\smallskip
\noindent\textbf{\texttt{\small DERIV:\allowbreak{}EARLIER\_\allowbreak{}PROOF\_\allowbreak{}RECOVERY:\allowbreak{}PHYSICAL\_\allowbreak{}CARRIER\_\allowbreak{}REALIZATION}} \hfill {\footnotesize open}\par
{\small Construct actual cutoff sources, residue floor, collapsing isolated intervals, true-sector gap, physical clock and limiting identifications needed to apply the positive matrix atom theorem.}\par
{\footnotesize \srcpath{docs/derivations/earlier-proof-recovery.md} --- \#\# E18. Explicit remaining application obligations}\par
{\footnotesize \textit{Hypotheses.} The exact source assumptions and required estimates described in E18.}\par
{\footnotesize \textit{Remaining.} Explicit pending application or extension; neither the exact finite controls nor the narrower recovered theorem discharge it.}\par
\smallskip
\noindent\textbf{\texttt{\small DERIV:\allowbreak{}EARLIER\_\allowbreak{}PROOF\_\allowbreak{}RECOVERY:\allowbreak{}STAPLE\_\allowbreak{}TRANSVERSALITY\_\allowbreak{}REPAIR}} \hfill {\footnotesize open}\par
{\small Repair the force-rank and tube estimates with the full moving-adjoint derivative.}\par
{\footnotesize \srcpath{docs/derivations/earlier-proof-recovery.md} --- \#\# E18. Explicit remaining application obligations}\par
{\footnotesize \textit{Hypotheses.} The exact source assumptions and required estimates described in E18.}\par
{\footnotesize \textit{Remaining.} Explicit pending application or extension; neither the exact finite controls nor the narrower recovered theorem discharge it.}\par
\smallskip
\noindent\textbf{\texttt{\small DERIV:\allowbreak{}EARLIER\_\allowbreak{}PROOF\_\allowbreak{}RECOVERY:\allowbreak{}VSU\_\allowbreak{}WHOLE\_\allowbreak{}SPACE}} \hfill {\footnotesize open}\par
{\small Establish compatible normalization and uniform local estimates for the whole-space VSU limit, then justify the claimed far-field law.}\par
{\footnotesize \srcpath{docs/derivations/earlier-proof-recovery.md} --- \#\# E18. Explicit remaining application obligations}\par
{\footnotesize \textit{Hypotheses.} The exact source assumptions and required estimates described in E18.}\par
{\footnotesize \textit{Remaining.} Explicit pending application or extension; neither the exact finite controls nor the narrower recovered theorem discharge it.}\par
\smallskip
\noindent\textbf{\texttt{\small DERIV:\allowbreak{}UNIVERSAL\_\allowbreak{}CELLULAR\_\allowbreak{}HODGE:\allowbreak{}PROPER\_\allowbreak{}RETURN\_\allowbreak{}Q\_\allowbreak{}ANNIHILATION}} \hfill {\footnotesize open}\par
{\small The 60 paths flagged by the submitted endpoint-or-zero flux predicate identify physical tetrahedral retained-sector returns that vanish individually under intermediate Q projection.}\par
{\footnotesize \srcpath{docs/derivations/universal-cellular-hodge-tetrahedral.md} --- \#\#\# 3.4 Open physical-history identification (lines 133-156)}\par
{\footnotesize \textit{Hypotheses.} An actual Haar/Fierz Hilbert space, physical retained projector, insertion channels, and projected resolvent states matching the enumerated paths.}\par
{\footnotesize \textit{Remaining.} Open: construct and identify the physical states/projector and calculate each projected contribution. The 60/36 flux count and Q psi=0 do not establish this identification.}\par
\smallskip
\noindent\textbf{\texttt{\small DERIV:\allowbreak{}W6\_\allowbreak{}CONDITIONAL\_\allowbreak{}SCORE\_\allowbreak{}TAIL\_\allowbreak{}CONTROL:\allowbreak{}SCORE\_\allowbreak{}DOMINATION\_\allowbreak{}M10}} \hfill {\footnotesize open}\par
{\small For the actual compact square and specified source-preserving dilation, prove K\_\allowbreak{}g(w)<=C0+C1 W\_\allowbreak{}g(w) almost everywhere in nu\_\allowbreak{}g, where W\_\allowbreak{}g(w)=g\textasciicircum{}-2 E\_\allowbreak{}(mu\_\allowbreak{}g)(V|w), with finite nonnegative C0,C1 independent of 0<g<g\_\allowbreak{}*. This sufficient conditional-score domination is open.}\par
{\footnotesize \srcpath{docs/derivations/w6-conditional-score-tail-control.md} --- \#\#\# M.4. The precise remaining conditional-score inequality (lines 473-495)}\par
{\footnotesize \textit{Hypotheses.} Actual fixed twelve-edge, four-face compact Wilson square, true normalized quantum ground Psi\_\allowbreak{}g, dmu\_\allowbreak{}g=Psi\_\allowbreak{}g\textasciicircum{}2 dU, literal source w=Tr(U2 U3)/2 and its marginal nu\_\allowbreak{}g.; sigma\_\allowbreak{}g=(partial\_\allowbreak{}g Psi\_\allowbreak{}g+D Psi\_\allowbreak{}g/g)/Psi\_\allowbreak{}g uses the specified compact dilation of Q8, including its Haar divergence; K\_\allowbreak{}g=Var\_\allowbreak{}(mu\_\allowbreak{}g)(sigma\_\allowbreak{}g|w) retains all compact corrections and rare coarse fibers.; The established actual fixed-square gap and ground-energy bounds hold on the specified small-coupling interval. Averaged conditional-score control is not assumed to imply pointwise domination.}\par
{\footnotesize \textit{Remaining.} Prove M10 for the actual conditional law, or retain its open status if a different sufficient source-energy estimate is proved. The phase-adapted local implication requires whole-fiber tails, conditional-mean gluing and consistent global dilation before it can address this target.}\par
\smallskip
\noindent\textbf{\texttt{\small DERIV:\allowbreak{}W6\_\allowbreak{}CONDITIONAL\_\allowbreak{}TRANSPORT\_\allowbreak{}OBSTRUCTION:\allowbreak{}M10\_\allowbreak{}SYNCHRONIZED\_\allowbreak{}SCORE}} \hfill {\footnotesize open}\par
{\small The selected synchronized realization of the generic M10 target asks whether the actual score for the specified S13 compact field obeys K\_\allowbreak{}g(w)<=C0+C1 g\textasciicircum{}-2 E\_\allowbreak{}mu\_\allowbreak{}g[V|w] with finite nonnegative constants uniform in sufficiently small positive g and almost every retained trace w.}\par
{\footnotesize \srcpath{docs/derivations/w6-conditional-transport-obstruction.md} --- \#\# S6 Remaining all-fiber estimate (lines 317-375), explicitly open synchronized M10 target}\par
{\footnotesize \textit{Hypotheses.} Actual positive four-face SU(2) square ground and literal retained trace w=tr(U2 U3)/2.; The particular synchronized four radial profiles S13 and their full Haar unitary generator, not an arbitrary Q8 field.; All conditional source regimes, including rare fibers and the antipodal degeneration, are retained in the requested uniform estimate.}\par
{\footnotesize \textit{Remaining.} Prove the actual full conditional score bound, or a sufficient direct source-weighted pairing estimate. The original M11-M15 consequence then becomes applicable;}\par
\smallskip
\noindent\textbf{\texttt{\small DERIV:\allowbreak{}W6\_\allowbreak{}GROUND\_\allowbreak{}JETS\_\allowbreak{}TRANSPORT\_\allowbreak{}BUDGET:\allowbreak{}INTERACTING\_\allowbreak{}GRID\_\allowbreak{}COMPARISON}} \hfill {\footnotesize open}\par
{\small For the actual interacting growing-grid family, obtain a volume-independent fast floor and complete source/transport energy estimates using linked or locally centered energy bounds, and control the coarse-theory matching discrepancy after the prescribed injection. The fixed-square constants depending on total ground energy E and the Gaussian grid floor do not prove this extension.}\par
{\footnotesize \srcpath{docs/derivations/w6-ground-jets-and-transport-budget.md} --- \#\# 6. Source energies, subdivision, interacting grids, and the physical clock (lines 351-366)}\par
{\footnotesize \textit{Hypotheses.} Actual coupled Wilson grids and specified block/source injections, with true vacuum subtraction, common physical energy normalization and the required retained, fast and spectral-energy domains.; Constants must be independent of growing spatial volume in the intended small-coupling regime; the existing fixed-square E and gamma are not assumed to have that uniformity.; Actual source-energy comparison factors and the discrepancy between the generated coarse form and prescribed coarse theory are retained; the Gaussian covariance-weighted source is not identified with the nonlinear quantum source without proof.}\par
{\footnotesize \textit{Remaining.} Prove the interacting fast-floor and connected energy/source comparison in the specified grid family. Controlling the local coupling remainder does not control coarse matching, the explicit physical-clock change in R14, or convergence of the exact transported jets and observables.}\par
\smallskip
\noindent\textbf{\texttt{\small DERIV:\allowbreak{}W6\_\allowbreak{}GROUND\_\allowbreak{}JETS\_\allowbreak{}TRANSPORT\_\allowbreak{}BUDGET:\allowbreak{}SOURCE\_\allowbreak{}ENERGY\_\allowbreak{}JETS\_\allowbreak{}R10}} \hfill {\footnotesize open}\par
{\small For the actual complete compact-square source/vacuum generator A(g) in R9, prove ||A\textasciicircum{}(r)(g)||\_\allowbreak{}(q\_\allowbreak{}g->q\_\allowbreak{}g)<=d\_\allowbreak{}r g\textasciicircum{}(-r-1) for r=0,1,2 with finite constants independent of small positive g. This is the open model input R10, not the already proved conditional implication R10 to R12.}\par
{\footnotesize \srcpath{docs/derivations/w6-ground-jets-and-transport-budget.md} --- \#\# 4. An explicit bridge from the actual transport generator to M1,M2,M3 (lines 199-218); section 5 (lines 271-309)}\par
{\footnotesize \textit{Hypotheses.} Actual fixed compact-square Hamiltonian and common electric H1 form domain; actual simple normalized positive ground, 0<=e\_\allowbreak{}g<=E and physical gap gamma>0 on 0<g<g\_\allowbreak{}*.; P\_\allowbreak{}g is the complete true-ground literal-source projection from C1, and A(g) is exactly the R9 source/vacuum generator; q\_\allowbreak{}g[u]=H\_\allowbreak{}g[u]+gamma||u||\textasciicircum{}2, including vacuum components.; Parameter derivatives are taken at positive g on the fixed compact form domain. The target covers all source and fast energies, including spatial derivatives and generator derivatives through order two.}\par
{\footnotesize \textit{Remaining.} Establish the conditional-source projection energy bounds beyond the proved vacuum-only generator. A direct proof of mixed-form R11 is an alternative sufficient route and need not prove R10 itself.}\par
\smallskip
\noindent\textbf{\texttt{\small DERIV:\allowbreak{}W6\_\allowbreak{}SOURCE\_\allowbreak{}ENERGY\_\allowbreak{}JETS\_\allowbreak{}R10:\allowbreak{}H0\_\allowbreak{}ENERGY\_\allowbreak{}CONTRACTION}} \hfill {\footnotesize open}\par
{\small Establish the proposed bound b\_\allowbreak{}g[Pi\_\allowbreak{}g Phi] <= q\_\allowbreak{}g[Omega\_\allowbreak{}g Phi] with kappa\_\allowbreak{}P=1 on the actual model; the weaker uniform finite-kappa\_\allowbreak{}P H0 also suffices for R10.}\par
{\footnotesize \srcpath{docs/derivations/w6-source-energy-jets-r10.md} --- \#\# 6. Scaling defects and unresolved estimates (C6)}\par
{\footnotesize \textit{Hypotheses.} Actual fixed compact square at positive g; normalized positive ground, q\_\allowbreak{}g=H\_\allowbreak{}g+gamma I, common form domain.; Constants must be independent of g on 0<g<g\_\allowbreak{}*.}\par
{\footnotesize \textit{Remaining.} The submitted proof is not a discharge. C1-C5 repair identities and conditional reductions;}\par
\smallskip
\noindent\textbf{\texttt{\small DERIV:\allowbreak{}W6\_\allowbreak{}SOURCE\_\allowbreak{}ENERGY\_\allowbreak{}JETS\_\allowbreak{}R10:\allowbreak{}H1\_\allowbreak{}FIBERWISE\_\allowbreak{}SCORE\_\allowbreak{}VARIANCE}} \hfill {\footnotesize open}\par
{\small Establish sup\_\allowbreak{}w Var\_\allowbreak{}mu\_\allowbreak{}g(sigma\_\allowbreak{}g|w) <= kappa\_\allowbreak{}0 g\textasciicircum{}-2 uniformly at small positive coupling.}\par
{\footnotesize \srcpath{docs/derivations/w6-source-energy-jets-r10.md} --- \#\# 6. Scaling defects and unresolved estimates (C6)}\par
{\footnotesize \textit{Hypotheses.} Actual fixed compact square at positive g; normalized positive ground, q\_\allowbreak{}g=H\_\allowbreak{}g+gamma I, common form domain.; Constants must be independent of g on 0<g<g\_\allowbreak{}*.}\par
{\footnotesize \textit{Remaining.} The submitted proof is not a discharge. C1-C5 repair identities and conditional reductions;}\par
\smallskip
\noindent\textbf{\texttt{\small DERIV:\allowbreak{}W6\_\allowbreak{}SOURCE\_\allowbreak{}ENERGY\_\allowbreak{}JETS\_\allowbreak{}R10:\allowbreak{}KATO\_\allowbreak{}DERIVATIVES}} \hfill {\footnotesize open}\par
{\small Establish ||partial\_\allowbreak{}g\textasciicircum{}r[P\_\allowbreak{}g prime,P\_\allowbreak{}g]||\_\allowbreak{}(q\_\allowbreak{}g->q\_\allowbreak{}g) <= d\_\allowbreak{}P,r g\textasciicircum{}(-r-1), r=0,1,2, for the actual source projection.}\par
{\footnotesize \srcpath{docs/derivations/w6-source-energy-jets-r10.md} --- \#\# 6. Scaling defects and unresolved estimates (C6)}\par
{\footnotesize \textit{Hypotheses.} Actual fixed compact square at positive g; normalized positive ground, q\_\allowbreak{}g=H\_\allowbreak{}g+gamma I, common form domain.; Constants must be independent of g on 0<g<g\_\allowbreak{}*.}\par
{\footnotesize \textit{Remaining.} The submitted proof is not a discharge. C1-C5 repair identities and conditional reductions;}\par
\smallskip
\noindent\textbf{\texttt{\small DERIV:\allowbreak{}W6\_\allowbreak{}SOURCE\_\allowbreak{}ENERGY\_\allowbreak{}JETS\_\allowbreak{}R10:\allowbreak{}VACUUM\_\allowbreak{}CROSS\_\allowbreak{}DERIVATIVES}} \hfill {\footnotesize open}\par
{\small Establish ||partial\_\allowbreak{}g\textasciicircum{}r A\_\allowbreak{}nu||\_\allowbreak{}(q\_\allowbreak{}g->q\_\allowbreak{}g) <= d\_\allowbreak{}nu,r g\textasciicircum{}(-r-1), r=0,1,2, for the actual vacuum-cross source vector.}\par
{\footnotesize \srcpath{docs/derivations/w6-source-energy-jets-r10.md} --- \#\# 6. Scaling defects and unresolved estimates (C6)}\par
{\footnotesize \textit{Hypotheses.} Actual fixed compact square at positive g; normalized positive ground, q\_\allowbreak{}g=H\_\allowbreak{}g+gamma I, common form domain.; Constants must be independent of g on 0<g<g\_\allowbreak{}*.}\par
{\footnotesize \textit{Remaining.} The submitted proof is not a discharge. C1-C5 repair identities and conditional reductions;}\par
\smallskip
\noindent\textbf{\texttt{\small DERIV:\allowbreak{}W6\_\allowbreak{}SOURCE\_\allowbreak{}GENERATOR\_\allowbreak{}SCORE\_\allowbreak{}FRAME:\allowbreak{}CONDITIONAL\_\allowbreak{}ENERGY\_\allowbreak{}H0\_\allowbreak{}H4}} \hfill {\footnotesize open}\par
{\small Prove, uniformly on 0<g<g\_\allowbreak{}*, the fiberwise conditional-energy estimates H0-H4 for the undilated ground score on level sets of w. Their averaged forms are established (SF7b and the integrated J\_\allowbreak{}g bound conditional on H4); the fiberwise forms are open.}\par
{\footnotesize \srcpath{docs/derivations/w6-source-generator-score-frame.md} --- \#\# 5. Exact reduction of the order-zero R10 bound (SF8-SF9), conditional-energy hypotheses}\par
{\footnotesize \textit{Hypotheses.} Actual fixed square, undilated score partial\_\allowbreak{}g log Omega\_\allowbreak{}g; K\textasciicircum{}0\_\allowbreak{}g is not the dilated K\_\allowbreak{}g of M10.}\par
{\footnotesize \textit{Remaining.} A fiberwise Poincare-type inequality for -Delta\_\allowbreak{}mu restricted to level sets of w, uniform in g, would give H1 and H2 from SF6. H0, H3 and H4 concern the source energy of conditional expectations and need a separate argument.}\par
\smallskip
\noindent\textbf{\texttt{\small DERIV:\allowbreak{}W6\_\allowbreak{}SYNCHRONIZED\_\allowbreak{}M10\_\allowbreak{}DOMINATION:\allowbreak{}AMPLITUDE\_\allowbreak{}GRADIENT\_\allowbreak{}SCALING}} \hfill {\footnotesize open}\par
{\small The actual uniform amplitude-score gradient bound ca/g remains open; changing parameter to h=g\textasciicircum{}2 and a spatial WKB expansion do not alone prove it.}\par
{\footnotesize \srcpath{docs/derivations/w6-synchronized-m10-domination.md} --- \#\# 5. Correct parameter-derivative criterion (lines 144-171)}\par
{\footnotesize \textit{Hypotheses.} Actual positive ground amplitude in a valid phase tube.}\par
{\footnotesize \textit{Remaining.} Establish uniform spatial control of h*partial\_\allowbreak{}h log A and Z log A for the actual ground; the logarithmic-parameter criterion gives the sufficient implication.}\par
\smallskip
\noindent\textbf{\texttt{\small DERIV:\allowbreak{}W6\_\allowbreak{}SYNCHRONIZED\_\allowbreak{}M10\_\allowbreak{}DOMINATION:\allowbreak{}FULL\_\allowbreak{}SCORE\_\allowbreak{}DOMINATION\_\allowbreak{}M10}} \hfill {\footnotesize open}\par
{\small Actual full-fiber M10 for S13 remains open. The earlier completion claim is retracted; the repaired assembly implication retains all unsupplied model hypotheses.}\par
{\footnotesize \srcpath{docs/derivations/w6-synchronized-m10-domination.md} --- \#\# 7. Assembly criterion and actual status (lines 262-286)}\par
{\footnotesize \textit{Hypotheses.} Actual four-face compact SU(2) quantum ground law and S13 dilation.}\par
{\footnotesize \textit{Remaining.} Actual complement is discharged. Supply the polynomial relative-amplitude and tube moment estimates, intermediate controls, and actual normal/angular comparison and score error.}\par
\smallskip
\noindent\textbf{\texttt{\small DERIV:\allowbreak{}WILSON\_\allowbreak{}MARKED\_\allowbreak{}TRANSFER:\allowbreak{}WT6}} \hfill {\footnotesize open}\par
{\small The proposed rooted coefficient-operator bound WT6 is an explicit sufficient route to complete-shell transport; this historical argument did not establish it. Its continuation uses a different full-Hilbert estimate and finite-order locality.}\par
{\footnotesize \srcpath{docs/derivations/wilson-marked-transfer.md} --- \#\# WT-6. The remaining complete-shell estimate and conditional transport (lines 209-248)}\par
{\footnotesize \textit{Hypotheses.} Finite cubic link Hilbert product; positive free generators kill the product vacuum and have gap delta>0 on its complement.; Bounded plaquette interactions ||V\_\allowbreak{}p||<=v; physical character-clock step tau>0 and blocks of duration ell=m*tau.; Actual vacuum-dressed transfer on the specified exact-support coefficient space must first be constructed.}\par
{\footnotesize \textit{Remaining.} Formalize or discharge the particular WT6 rooted norm if used; do not equate the later sufficient route with a proof of this specific estimate.}\par
\smallskip
\noindent\textbf{\texttt{\small DERIV:\allowbreak{}WILSON\_\allowbreak{}SC17\_\allowbreak{}SPATIAL\_\allowbreak{}CLOSURE:\allowbreak{}ACTUAL\_\allowbreak{}REFERENCE\_\allowbreak{}DEFECT\_\allowbreak{}BOUND}} \hfill {\footnotesize open}\par
{\small On the specified actual Wilson conditional source chart, establish a volume-uniform weighted bound on the complete signed Hessian defect D\_\allowbreak{}B=Hess(V\_\allowbreak{}B+(H\_\allowbreak{}out Omega)/Omega)-2epsilon A\textasciicircum{}2, including the actual kinetic metric, moving source projector and localization terms, with norm(D\_\allowbreak{}B)\_\allowbreak{}(s,infty)<=delta and 2 beta\_\allowbreak{}s(A)\textasciicircum{}2 delta/epsilon<1.}\par
{\footnotesize \srcpath{docs/derivations/wilson-sc17-spatial-closure.md} --- \#\# 4. The actual quantum pressure and cutoff defects are part of the input (lines 698-756), with the R5-R7 budget in section 2}\par
{\footnotesize \textit{Hypotheses.} Actual positive Wilson vacuum, compatible source and block coordinates, and multiplication terms assigned exactly once; the complete coupled Gaussian reference is subtracted.; The reference semigroup budget beta\_\allowbreak{}s(A) is finite in the chosen weighted norm at the fixed blocking scale; epsilon>0 and all conditional, metric and cutoff domains are retained.; The desired uniformity concerns the stated block family and surrounding volume; a separate source-compatible reference-angle margin is still needed for global assembly.}\par
{\footnotesize \textit{Remaining.} Prove the actual complete weighted smallness estimate. The proved Gaussian pressure cancellation, an isolated magnetic Taylor coefficient, or formalizing the conditional Riccati implication does not discharge this model obligation.}\par
\smallskip
\noindent\textbf{\texttt{\small DERIV:\allowbreak{}WILSON\_\allowbreak{}SELECTED\_\allowbreak{}INVERSE\_\allowbreak{}WALL:\allowbreak{}W6}} \hfill {\footnotesize open}\par
{\small The needed actual interacting Wilson bound is sup\_\allowbreak{}(Lambda,b[p]=1)|d\_\allowbreak{}g,QQ[R0 t1p,Rg t1p]|<=C|g| uniformly in volume, coupling, backgrounds and required spectral interval.}\par
{\footnotesize \srcpath{docs/derivations/wilson-selected-inverse-wall.md} --- \#\# 5. The exact point at which this argument cannot proceed (lines 221-264)}\par
{\footnotesize \textit{Hypotheses.} The operator chart, vacuum subtraction and fast/retained spaces are those of the spatial Schur derivation; domains are explicitly checked rather than identified by notation.; 0<g<=g0; complete relevant retained energy space and high-retained contribution controlled.}\par
{\footnotesize \textit{Remaining.} Prove the actual coupled signed resolvent pairing with the complete dressed force, physical energy normalization and an interacting fast floor; no existing Gaussian control discharges it.}\par
\smallskip
\noindent\textbf{\texttt{\small DERIV:\allowbreak{}YANGMILLS\_\allowbreak{}CONTINUUM\_\allowbreak{}BALABAN\_\allowbreak{}MULTISCALE\_\allowbreak{}PROOF:\allowbreak{}OS1}} \hfill {\footnotesize open}\par
{\small Full Euclidean covariance requires restoration of SO(4) beyond cubic block symmetry; the repair record ties this to the unproved irrelevance estimate.}\par
{\footnotesize \srcpath{docs/derivations/yangmills-continuum-balaban-multiscale-proof.md} --- \#\#\# 8.2 OS1: Euclidean Covariance (lines 375-382)}\par
{\footnotesize \textit{Hypotheses.} Submitted SU(N), N>=2 continuum multiscale construction; read together with its in-document Repair record audited 2026-09-09.; The source title and summary claim completion, but its repair record retains explicit averaging/interpolation, irrelevance, physical-scale and scattering obligations. No machine certification is inferred from prose.}\par
{\footnotesize \textit{Remaining.} Prove rotational restoration for the actual continuum limit, then translation/rotation covariance of its Schwinger distributions.}\par
\smallskip
\noindent\textbf{\texttt{\small DERIV:\allowbreak{}YANGMILLS\_\allowbreak{}CONTINUUM\_\allowbreak{}BALABAN\_\allowbreak{}MULTISCALE\_\allowbreak{}PROOF:\allowbreak{}THEOREM\_\allowbreak{}7\_\allowbreak{}1}} \hfill {\footnotesize open}\par
{\small The submitted Cauchy theorem needs a summable cross-scale expectation increment. The printed O(1/k) bound alone is not summable; the repair record names the missing irrelevance factor a\_\allowbreak{}k\textasciicircum{}theta.}\par
{\footnotesize \srcpath{docs/derivations/yangmills-continuum-balaban-multiscale-proof.md} --- \#\#\# 7.1 Cauchy Sequence of Effective Actions (lines 317-340)}\par
{\footnotesize \textit{Hypotheses.} Submitted SU(N), N>=2 continuum multiscale construction; read together with its in-document Repair record audited 2026-09-09.; The source title and summary claim completion, but its repair record retains explicit averaging/interpolation, irrelevance, physical-scale and scattering obligations. No machine certification is inferred from prose.; Actual gauge-invariant cylinder expectations and uniform operator normalization; a positive irrelevance exponent or another summable increment bound.}\par
{\footnotesize \textit{Remaining.} Prove the actual summable remainder after marginal matching, then construct the Cauchy limit with its observable norms and scale maps.}\par
\smallskip
\noindent\textbf{\texttt{\small DERIV:\allowbreak{}YANGMILLS\_\allowbreak{}CONTINUUM\_\allowbreak{}BALABAN\_\allowbreak{}MULTISCALE\_\allowbreak{}PROOF:\allowbreak{}THEOREM\_\allowbreak{}8\_\allowbreak{}2}} \hfill {\footnotesize open}\par
{\small The claimed continuum exponential clustering requires a strictly positive rate in common physical units at a fixed physical scale; UV cell-scale polymer decay alone does not provide that conclusion.}\par
{\footnotesize \srcpath{docs/derivations/yangmills-continuum-balaban-multiscale-proof.md} --- \#\#\# 8.5 OS4: Exponential Clustering and Mass Gap (lines 412-436)}\par
{\footnotesize \textit{Hypotheses.} Submitted SU(N), N>=2 continuum multiscale construction; read together with its in-document Repair record audited 2026-09-09.; The source title and summary claim completion, but its repair record retains explicit averaging/interpolation, irrelevance, physical-scale and scattering obligations. No machine certification is inferred from prose.; The repair record explicitly requires the missing bridge to a scale with g=O(1); uniform cutoff/volume estimates are retained.}\par
{\footnotesize \textit{Remaining.} Prove the actual nonperturbative KP or equivalent correlation estimate at a fixed physical reference scale, then transport it through the constructed continuum limit.}\par
\smallskip
\noindent\textbf{\texttt{\small DERIV:\allowbreak{}YANGMILLS\_\allowbreak{}CONTINUUM\_\allowbreak{}BALABAN\_\allowbreak{}MULTISCALE\_\allowbreak{}PROOF:\allowbreak{}GAUGE\_\allowbreak{}HESSIAN}} \hfill {\footnotesize disputed}\par
{\small The submission gives the background gauge Hessian (5.3)-(5.5); the operator must include the correct covariant 1-form curvature term and the repaired averaging map consistently.}\par
{\footnotesize \srcpath{docs/derivations/yangmills-continuum-balaban-multiscale-proof.md} --- \#\#\# 5.1 Fluctuation Gauge Fixing (lines 183-197)}\par
{\footnotesize \textit{Hypotheses.} Submitted SU(N), N>=2 continuum multiscale construction; read together with its in-document Repair record audited 2026-09-09.; The source title and summary claim completion, but its repair record retains explicit averaging/interpolation, irrelevance, physical-scale and scattering obligations. No machine certification is inferred from prose.; Smooth background and gauge-fixed fluctuation domain; repaired deterministic block constraint.}\par
{\footnotesize \textit{Remaining.} Formalize the covariant exterior-derivative/Weitzenbock identity with signs and domains; reconcile Q\_\allowbreak{}k versus Q\_\allowbreak{}B before using the Hessian.}\par
\smallskip
\noindent\textbf{\texttt{\small DERIV:\allowbreak{}YANGMILLS\_\allowbreak{}CONTINUUM\_\allowbreak{}BALABAN\_\allowbreak{}MULTISCALE\_\allowbreak{}PROOF:\allowbreak{}PROPOSITION\_\allowbreak{}3\_\allowbreak{}4}} \hfill {\footnotesize disputed}\par
{\small Proposition 3.4 asserts a unique gauge-equivariant constrained background minimizer with the Sobolev/coercivity bounds (3.9)-(3.12); the repair record says the corrected affine adjoint constraint and transported interpolation have not been propagated.}\par
{\footnotesize \srcpath{docs/derivations/yangmills-continuum-balaban-multiscale-proof.md} --- \#\#\# 3.2 The Background Field Variational Problem (lines 111-139)}\par
{\footnotesize \textit{Hypotheses.} Submitted SU(N), N>=2 continuum multiscale construction; read together with its in-document Repair record audited 2026-09-09.; The source title and summary claim completion, but its repair record retains explicit averaging/interpolation, irrelevance, physical-scale and scattering obligations. No machine certification is inferred from prose.; Small curvature g\_\allowbreak{}(k+1)\textasciicircum{}kappa a\_\allowbreak{}(k+1)\textasciicircum{}-2, 0<kappa<1/6, sufficiently small coupling; all correct gauge and boundary domains still required.}\par
{\footnotesize \textit{Remaining.} First implement the actual repaired deterministic constraint and equivariant interpolation, then prove coercivity, existence, uniqueness and smooth parameter dependence on its genuine domain.}\par
\smallskip
\noindent\textbf{\texttt{\small DERIV:\allowbreak{}YANGMILLS\_\allowbreak{}CONTINUUM\_\allowbreak{}BALABAN\_\allowbreak{}MULTISCALE\_\allowbreak{}PROOF:\allowbreak{}THEOREM\_\allowbreak{}5\_\allowbreak{}1}} \hfill {\footnotesize disputed}\par
{\small Theorem 5.1 asserts a uniform positive background fluctuation floor and exponential pointwise resolvent kernel bound. The actual averaging conjugation, covariant Poincare and kernel regularity estimates remain required.}\par
{\footnotesize \srcpath{docs/derivations/yangmills-continuum-balaban-multiscale-proof.md} --- \#\#\# 5.3 Combes--Thomas Resolvent Localization for Continuous \$\textbackslash{}mathcal\{H\}\_\allowbreak{}\{B\_\allowbreak{}k\}\$ (lines 209-248)}\par
{\footnotesize \textit{Hypotheses.} Submitted SU(N), N>=2 continuum multiscale construction; read together with its in-document Repair record audited 2026-09-09.; The source title and summary claim completion, but its repair record retains explicit averaging/interpolation, irrelevance, physical-scale and scattering obligations. No machine certification is inferred from prose.; delta<=1/(4224ell\textasciicircum{}2); alpha0>=8delta+1; the domain, cutoff and averaging operator must be specified.}\par
{\footnotesize \textit{Remaining.} Formalize the corrected covariant averaged operator and exponential conjugation including its nonlocal averaging term; replace or justify the all-x,y bounded continuum resolvent kernel claim, which has a diagonal singularity without an explicit smoothing cutoff.}\par
\smallskip
\noindent\textbf{\texttt{\small DERIV:\allowbreak{}YANGMILLS\_\allowbreak{}RESOLVENT\_\allowbreak{}LOCALIZATION\_\allowbreak{}AND\_\allowbreak{}DIRICHLET\_\allowbreak{}GAP:\allowbreak{}SUBMITTED\_\allowbreak{}W6\_\allowbreak{}CLOSURE}} \hfill {\footnotesize disputed}\par
{\small The proposal asserts actual Wilson W6 closure from massive spatial Green decay; its own source audit records that the computed massive scalar operator is not identified with Wilson.}\par
{\footnotesize \srcpath{docs/derivations/yangmills-resolvent-localization-and-dirichlet-gap.md} --- \#\# 1. Summary of Results (lines 20-39)}\par
{\footnotesize \textit{Hypotheses.} This submitted proposal is explicitly preserved under its September 9 source audit; closure claims are not the current graph verdict.; The actual model/operator identification is an obligation, separate from a scalar finite replay.}\par
{\footnotesize \textit{Remaining.} Establish the actual Wilson fast floor, connected signed pairing and physical source identification; the submission does not discharge these.}\par
\smallskip
\noindent\textbf{\texttt{\small DERIV:\allowbreak{}EARLIER\_\allowbreak{}PROOF\_\allowbreak{}RECOVERY:\allowbreak{}B6\_\allowbreak{}RETAINED\_\allowbreak{}KRYLOV\_\allowbreak{}CAMPAIGN}} \hfill {\footnotesize conditional}\par
{\small The retained SU(3) B6 certificate reports radial dimensions 67,155,133 and maximum full-space Ritz residual 8.775675697349887e-14 on 201 g-values in [1,2], with crossing g=1.3039546641713.}\par
{\footnotesize \srcpath{docs/derivations/earlier-proof-recovery.md} --- \#\# E6b. Retained B6 numerical campaign}\par
{\footnotesize \textit{Hypotheses.} Historical source K6(u)=E-uM, source coupling conversion and branch choices; level-eight 511-column word spaces; relative SVD cutoff 1e-10.; The certificate is retained and hash-verified; the B6 eigensolver was not rerun in this integration.}\par
{\footnotesize \textit{Remaining.} The original numerical certificate is preserved. A fresh B6 eigensolver replay and any uniform-in-coupling or full-shell conclusion are not supplied.}\par
\smallskip
\noindent\textbf{\texttt{\small DERIV:\allowbreak{}W6\_\allowbreak{}CONDITIONAL\_\allowbreak{}SCORE\_\allowbreak{}TAIL\_\allowbreak{}CONTROL:\allowbreak{}SCORE\_\allowbreak{}WEIGHT\_\allowbreak{}FROM\_\allowbreak{}M10}} \hfill {\footnotesize conditional}\par
{\small Conditional on the actual M10 domination, every centered finite-energy source satisfies integral K\_\allowbreak{}g|f|\textasciicircum{}2 dnu\_\allowbreak{}g <= [C1+(C0+C1 E)/gamma] b\_\allowbreak{}g[f], and the actual median Hardy constant obeys mathfrak\_\allowbreak{}B\_\allowbreak{}g<=C1+2(C0+C1 E)/gamma uniformly at small coupling. The implication is analytically proved; its M10 model input remains open.}\par
{\footnotesize \srcpath{docs/derivations/w6-conditional-score-tail-control.md} --- \#\#\# M.4. The precise remaining conditional-score inequality and M.5 (lines 497-563)}\par
{\footnotesize \textit{Hypotheses.} Actual compact-square source form b\_\allowbreak{}g[f]=g\textasciicircum{}2 integral (1-w\textasciicircum{}2)|f'|\textasciicircum{}2 dnu\_\allowbreak{}g, actual ground-energy bound e\_\allowbreak{}g<=E and physical gap gamma>0, with the true-ground source identity M7.; The actual conditional-score domination M10 holds with finite nonnegative C0,C1 uniform in g.; Sources belong to the full compact retained form domain; centered sources give M12, and the median-anchored half-source argument retains the uncentered vacuum-energy term for M14-M15.}\par
{\footnotesize \textit{Remaining.} Formalize this already established conditional implication if needed. Actual application requires M10 or another sufficient score estimate;}\par
\smallskip
\noindent\textbf{\texttt{\small DERIV:\allowbreak{}W6\_\allowbreak{}CONDITIONAL\_\allowbreak{}TRANSPORT\_\allowbreak{}OBSTRUCTION:\allowbreak{}LOCAL\_\allowbreak{}GLOBAL\_\allowbreak{}SCORE\_\allowbreak{}CRITERION}} \hfill {\footnotesize conditional}\par
{\small The stated uniform conditional tube moments, differentiated amplitude-score bound and synchronized phase remainder imply the local S14 score estimate. S15 glues it to the full conditional variance using the actual outside-tube squared deviation from the same reference beta\_\allowbreak{}g. Uniform potential-weighted complement and remaining-regime bounds would supply synchronized M10.}\par
{\footnotesize \srcpath{docs/derivations/w6-conditional-transport-obstruction.md} --- \#\# S6 Remaining all-fiber estimate (lines 317-375)}\par
{\footnotesize \textit{Hypotheses.} Conditional tube moments E|eta|\textasciicircum{}(2j)<=c\_\allowbreak{}(2j)g\textasciicircum{}(2j), j=1,2,3, with uniform constants.; True relative amplitude score has fast derivative at most c\_\allowbreak{}a/g; the phase is smooth with bounded third fast derivatives and the stated oscillator quadratic jet.; Synchronized tangency holds. The same reference beta\_\allowbreak{}g is used inside and outside the tube, so conditional mean differences remain in the outside deviation.; Full M10 additionally needs a uniform potential-weighted outside deviation bound and a degeneration-uniform actual score bound in the remaining source regimes.}\par
{\footnotesize \textit{Remaining.} Establish the actual uniform differentiated amplitude, tube-moment, conditional complement and antipodal estimates. Undifferentiated local comparison used for the obstruction does not discharge these hypotheses.}\par
\smallskip
\noindent\textbf{\texttt{\small DERIV:\allowbreak{}W6\_\allowbreak{}GROUND\_\allowbreak{}JETS\_\allowbreak{}TRANSPORT\_\allowbreak{}BUDGET:\allowbreak{}COMPLETE\_\allowbreak{}BUDGET}} \hfill {\footnotesize conditional}\par
{\small For the complete actual positive-reference source/vacuum transport, energy-operator bounds ||A\textasciicircum{}(r)(g)||\_\allowbreak{}(q\_\allowbreak{}g->q\_\allowbreak{}g)<=d\_\allowbreak{}r g\textasciicircum{}(-r-1), r=0,1,2, imply explicit all-energy M\_\allowbreak{}j(s)<=c\_\allowbreak{}j s\textasciicircum{}-j for j=1,2,3 and hence the complete compact Schur residual budget. The composite budget retains changing source energies and physical clock factors.}\par
{\footnotesize \srcpath{docs/derivations/w6-ground-jets-and-transport-budget.md} --- \#\# 4. An explicit bridge from the actual transport generator to M1,M2,M3 (lines 199-367)}\par
{\footnotesize \textit{Hypotheses.} Complete actual common-domain source/vacuum transport R'=A R and nonreducing graph from the positive-reference source theorem.; For q\_\allowbreak{}g=H\_\allowbreak{}g+gamma, full energy-space ||A\textasciicircum{}(r)(g)||<=d\_\allowbreak{}r g\textasciicircum{}(-r-1), r=0,1,2, with finite uniform constants.; Local t in [s/2,s], fixed spectral window below the actual gap. The composite source growth and subdivision inputs are the prior S2-S5 theorem.; Positive physical clock tau\_\allowbreak{}n; physical matching discrepancies are separately retained, not assumed zero.}\par
{\footnotesize \textit{Remaining.} Prove the actual complete conditional source-generator energy jets or exact mixed-form substitute; supply interacting-grid, coarse-matching and physical-clock convergence inputs.}\par
\smallskip
\noindent\textbf{\texttt{\small DERIV:\allowbreak{}W6\_\allowbreak{}SOURCE\_\allowbreak{}GENERATOR\_\allowbreak{}SCORE\_\allowbreak{}FRAME:\allowbreak{}R10\_\allowbreak{}ORDER\_\allowbreak{}ZERO\_\allowbreak{}REDUCTION}} \hfill {\footnotesize conditional}\par
{\small Under the conditional-energy hypotheses H0-H4 (uniform q\_\allowbreak{}g-boundedness of P\_\allowbreak{}g, sup\_\allowbreak{}w K\textasciicircum{}0\_\allowbreak{}g <= kappa\_\allowbreak{}0 g\textasciicircum{}-2, J\_\allowbreak{}g <= kappa\_\allowbreak{}1 g\textasciicircum{}-2 (1+g\textasciicircum{}-2 E[V|w]), the fast-to-source conditional energy bound, and b\_\allowbreak{}g[m\_\allowbreak{}g] <= kappa\_\allowbreak{}3 g\textasciicircum{}-2), the complete R9 generator satisfies ||A(g)||\_\allowbreak{}(q\_\allowbreak{}g->q\_\allowbreak{}g) <= d\_\allowbreak{}0 g\textasciicircum{}-1 with the\,\ldots}\par
{\footnotesize \srcpath{docs/derivations/w6-source-generator-score-frame.md} --- \#\# 5. Exact reduction of the order-zero R10 bound (SF8-SF9)}\par
{\footnotesize \textit{Hypotheses.} H0-H4 with constants independent of g on 0<g<g\_\allowbreak{}*; each is stated explicitly in section 5.; The rank-one q\_\allowbreak{}g operator bound and R4 from the transport-budget note; 0<=e\_\allowbreak{}g<=E and gap gamma>0.}\par
{\footnotesize \textit{Remaining.} Prove H0-H4. The r=1,2 cases of R10 are not reduced;}\par
\smallskip
\noindent\textbf{\texttt{\small DERIV:\allowbreak{}WILSON\_\allowbreak{}BACKGROUND\_\allowbreak{}COVARIANCE\_\allowbreak{}CONTINUATION:\allowbreak{}BF1\_\allowbreak{}BF3}} \hfill {\footnotesize conditional}\par
{\small A common fast form with relative defect norm theta\_\allowbreak{}z<1 gives the energy-weighted inverse difference <=theta\_\allowbreak{}z/(1-theta\_\allowbreak{}z), the selected signed pairing bound BF1, and the exact energy-dependent Schur/source identity BF3.}\par
{\footnotesize \srcpath{docs/derivations/wilson-background-covariance-continuation.md} --- \#\# BF1. The energy-dependent relative-form statement (lines 48-102)}\par
{\footnotesize \textit{Hypotheses.} Finite periodic cubic SU(2) lattice with nondegenerate plaquettes and side L >= 3.; Product unit-S3 electric metric; epsilon > 0, k >= 0, lambda = k/epsilon; actual positive normalized Wilson ground Omega.; A0(z)=F0-z>=f\_\allowbreak{}z I>0; V is a bounded A0-relative form with norm <=theta\_\allowbreak{}z<1.}\par
{\footnotesize \textit{Remaining.} The normalized bounded inverse, selected pairing and Schur minimization are formalized. Identify the actual common closed form, inverse-square-root coordinates and compatible cross-functionals for the Wilson operator.}\par
\smallskip
\noindent\textbf{\texttt{\small DERIV:\allowbreak{}WILSON\_\allowbreak{}BACKGROUND\_\allowbreak{}COVARIANCE\_\allowbreak{}CONTINUATION:\allowbreak{}SC17}} \hfill {\footnotesize conditional}\par
{\small The exact distant mixed derivative is -R\_\allowbreak{}xy(F\_\allowbreak{}e tensor V\_\allowbreak{}f+V\_\allowbreak{}e tensor F\_\allowbreak{}f)-F\_\allowbreak{}e tensor F\_\allowbreak{}f; the original all-coupling proof does not supply a summable distance envelope.}\par
{\footnotesize \srcpath{docs/derivations/wilson-background-covariance-continuation.md} --- \#\# 5. The neutral channel and the exact remaining spatial bound (lines 494-534)}\par
{\footnotesize \textit{Hypotheses.} Finite periodic cubic SU(2) lattice with nondegenerate plaquettes and side L >= 3.; Product unit-S3 electric metric; epsilon > 0, k >= 0, lambda = k/epsilon; actual positive normalized Wilson ground Omega.; Distinct distant tails; V\_\allowbreak{}ef=0; R\_\allowbreak{}xy is the actual charged inverse.}\par
{\footnotesize \textit{Remaining.} Formalize this identity and the singlet obstruction; the all-coupling spatially summable estimate still needs its complete quantum reference-defect bound.}\par
\smallskip
\noindent\textbf{\texttt{\small DERIV:\allowbreak{}WILSON\_\allowbreak{}MARKED\_\allowbreak{}SHELL\_\allowbreak{}TRANSPORT:\allowbreak{}S18\_\allowbreak{}S20}} \hfill {\footnotesize conditional}\par
{\small The common weighted Taylor majorant sums finite-order Wilson/Hamiltonian matching; the relative shell gap stays >(t3/4)u\textasciicircum{}2 q(k) away from Gamma and carrier projections converge by S20.}\par
{\footnotesize \srcpath{docs/derivations/wilson-marked-shell-transport.md} --- \#\# 6. Sum the matching and transport the carrier (lines 391-461)}\par
{\footnotesize \textit{Hypotheses.} SU(3), fixed spatial spacing and small magnetic coupling on the S1 common domain.; Calibrated kinetic-window and the actual September 5 Hamiltonian G18 creator/activity/complete-band hypotheses and literal source charts are explicit inputs.; The Hamiltonian G18 gap and kinetic-window coefficient comparison apply on the same weighted literal-source frame; epsilon is sufficiently small.}\par
{\footnotesize \textit{Remaining.} Formalize power-series tail convergence, the common polar source chart, relative spectral perturbation and exact symmetry cancellation at Gamma.}\par
\smallskip
\noindent\textbf{\texttt{\small DERIV:\allowbreak{}WILSON\_\allowbreak{}SC17\_\allowbreak{}SPATIAL\_\allowbreak{}CLOSURE:\allowbreak{}PROPOSED\_\allowbreak{}QUANTUM\_\allowbreak{}PARAMETER\_\allowbreak{}WINDOW}} \hfill {\footnotesize conditional}\par
{\small Under the proposed numerical reference and normalized-defect inputs, the scalar Riccati parameter satisfies p(lambda)<=p(3/100)=96228*sqrt(3)/625000<27/100<1, and the small root satisfies 0<=r\_\allowbreak{}minus<a\_\allowbreak{}min throughout 0<=lambda<=3/100. These are conditional parameter conclusions; the section's complete Wilson defect estimate and exponential spatial decay are not established by this arithmetic.}\par
{\footnotesize \srcpath{docs/derivations/wilson-sc17-spatial-closure.md} --- \#\# 6. Formal repair of the spatial closure for lambda >= 1/73 via the quantum reference defect (lines 781-821)}\par
{\footnotesize \textit{Hypotheses.} Real coupling 0<=lambda<=3/100; epsilon>0 when interpreting the normalized defect c=delta/(2epsilon).; The companion script's proposed defaults are L=3, C\_\allowbreak{}s=6/5 and reference velocity one; beta=(108/25)*sqrt(33), c=lambda*sqrt(lambda)/48 and a\_\allowbreak{}min=1/(3*sqrt(33)).; p(lambda)=4*beta\textasciicircum{}2*c and r\_\allowbreak{}minus=(1-sqrt(1-p(lambda)))/(2*beta).; These numerical inputs are supplied parameters. Their identification with a uniform semigroup budget, the complete conditional quantum-pressure defect and the actual reference floor is a separate model obligation, not a consequence of the quartic Taylor coefficient.}\par
{\footnotesize \textit{Remaining.} The scalar comparison is proved in the supported successor SC17\_\allowbreak{}P. To apply it to the source section's spatial-closure claim, establish the actual weighted reference budget and complete outside-pressure, metric, projector and cutoff defect bound, derive the ground-state Riccati comparison, and justify an exponential spatial weight.}\par
\smallskip
\noindent\textbf{\texttt{\small DERIV:\allowbreak{}WILSON\_\allowbreak{}SC17\_\allowbreak{}SPATIAL\_\allowbreak{}CLOSURE:\allowbreak{}R10\_\allowbreak{}R11}} \hfill {\footnotesize conditional}\par
{\small Uniform block Hessian floors 2a\_\allowbreak{}B,2a\_\allowbreak{}C and cross norm 2h\_\allowbreak{}BC yield c\_\allowbreak{}BC<=h\_\allowbreak{}BC/sqrt(a\_\allowbreak{}B a\_\allowbreak{}C); a reference angle budget plus full Hessian error r yields physical gap >=2epsilon[a0(1-kappa\_\allowbreak{}A)-2r].}\par
{\footnotesize \srcpath{docs/derivations/wilson-sc17-spatial-closure.md} --- \#\# 3. An explicit true-angle consequence if the full Hessian defect is known (lines 649-697)}\par
{\footnotesize \textit{Hypotheses.} Finite periodic cubic SU(2) lattice with nondegenerate plaquettes and side L >= 3.; Product unit-S3 electric metric; epsilon > 0, k >= 0, lambda = k/epsilon; actual positive normalized Wilson ground Omega.; a\_\allowbreak{}B,a\_\allowbreak{}C>0; h\_\allowbreak{}BC\textasciicircum{}2<a\_\allowbreak{}B a\_\allowbreak{}C; kappa\_\allowbreak{}A<1; 2r<a0(1-kappa\_\allowbreak{}A).}\par
{\footnotesize \textit{Remaining.} Formalize Brascamp-Lieb conditional variance, marginal Schur curvature and the actual projection-angle assembly.}\par
\smallskip
\noindent\textbf{\texttt{\small DERIV:\allowbreak{}WILSON\_\allowbreak{}SC17\_\allowbreak{}SPATIAL\_\allowbreak{}CLOSURE:\allowbreak{}R12\_\allowbreak{}R14}} \hfill {\footnotesize conditional}\par
{\small The necessary full conditional Hessian defect is Hess(V\_\allowbreak{}B+((H\_\allowbreak{}out Omega)/Omega))-2epsilon A\textasciicircum{}2; exact Gaussian pressure cancels the reference BB* term, while cutoffs create the explicit gradient and commutator defects.}\par
{\footnotesize \srcpath{docs/derivations/wilson-sc17-spatial-closure.md} --- \#\# 4. The actual quantum pressure and cutoff defects are part of the input (lines 698-756)}\par
{\footnotesize \textit{Hypotheses.} Finite periodic cubic SU(2) lattice with nondegenerate plaquettes and side L >= 3.; Product unit-S3 electric metric; epsilon > 0, k >= 0, lambda = k/epsilon; actual positive normalized Wilson ground Omega.; Actual vacuum and compatible Euclidean block chart; all multiplication terms assigned once.}\par
{\footnotesize \textit{Remaining.} Formalize the outside-pressure derivative, Gaussian cancellation and every cutoff/product-domain term; the actual uniform smallness estimate remains unproved.}\par
\smallskip
\noindent\textbf{\texttt{\small DERIV:\allowbreak{}WILSON\_\allowbreak{}SC17\_\allowbreak{}SPATIAL\_\allowbreak{}CLOSURE:\allowbreak{}R4\_\allowbreak{}R9}} \hfill {\footnotesize conditional}\par
{\small If beta\_\allowbreak{}s(A)<infinity and the full conditional potential Hessian defect is <=delta with 2 beta\_\allowbreak{}s(A)\textasciicircum{}2 delta/epsilon<1, the true log-ground Hessian error is <=r\_\allowbreak{}minus and conditional gap >=2epsilon(a-r\_\allowbreak{}minus).}\par
{\footnotesize \srcpath{docs/derivations/wilson-sc17-spatial-closure.md} --- \#\# 2. A rigorous conditional small-defect criterion (lines 574-648)}\par
{\footnotesize \textit{Hypotheses.} Finite periodic cubic SU(2) lattice with nondegenerate plaquettes and side L >= 3.; Product unit-S3 electric metric; epsilon > 0, k >= 0, lambda = k/epsilon; actual positive normalized Wilson ground Omega.; A>=aI>0; the full conditional potential includes outside quantum pressure; smooth ground convergence and domain conditions in the section hold.}\par
{\footnotesize \textit{Remaining.} Formalize the noncommutative weighted norm, Duhamel Riccati comparison, ground limit and curvature-to-gap passage.}\par
\smallskip
\noindent\textbf{\texttt{\small DERIV:\allowbreak{}WILSON\_\allowbreak{}SELECTED\_\allowbreak{}INVERSE\_\allowbreak{}WALL:\allowbreak{}W4\_\allowbreak{}W5}} \hfill {\footnotesize conditional}\par
{\small On admissible reference graph vectors, S\_\allowbreak{}g-S\_\allowbreak{}0=A\_\allowbreak{}g-r\_\allowbreak{}g*R\_\allowbreak{}g r\_\allowbreak{}g; only the selected inverse bound W5 plus direct/force remainders is sufficient for the displayed O(g\textasciicircum{}3) comparison.}\par
{\footnotesize \srcpath{docs/derivations/wilson-selected-inverse-wall.md} --- \#\# 3. Repair: only compare inverses on the graph force (lines 101-157)}\par
{\footnotesize \textit{Hypotheses.} The operator chart, vacuum subtraction and fast/retained spaces are those of the spatial Schur derivation; domains are explicitly checked rather than identified by notation.; z lies below both full fast spectra; J\_\allowbreak{}z p and the interacting minimizer are admissible; W5 includes f>0,a3,t2,t0,c\_\allowbreak{}sel.}\par
{\footnotesize \textit{Remaining.} The bounded Hilbert-space completion/minimum and weighted pairing are formalized. Construct the actual unbounded fast-form domain, source identification and reduced closed-form Schur operator.}\par
\smallskip
\noindent\textbf{\texttt{\small DERIV:\allowbreak{}WILSON\_\allowbreak{}SPATIAL\_\allowbreak{}SCHUR\_\allowbreak{}EXCESS:\allowbreak{}SP1\_\allowbreak{}SP6}} \hfill {\footnotesize conditional}\par
{\small The exact full-reference Schur excess is A\_\allowbreak{}g-V\_\allowbreak{}g*(I+C\_\allowbreak{}g)\textasciicircum{}-1 V\_\allowbreak{}g; the minimizing graph, two-sided form bounds and -dS\_\allowbreak{}g/dz=J\_\allowbreak{}g,z*J\_\allowbreak{}g,z retain the induced metric.}\par
{\footnotesize \srcpath{docs/derivations/wilson-spatial-schur-excess.md} --- \#\# 2. Exact excess about the full reference graph (lines 42-114)}\par
{\footnotesize \textit{Hypotheses.} Closed self-adjoint reference and interacting forms on the stated compatible P/Q domains, already expressed in a common physical source chart and vacuum-subtracted.; Real energy z below the full fast restriction; positive retained energy weight B.; F\_\allowbreak{}z>=f\_\allowbreak{}zI>0; ||C\_\allowbreak{}g||<=theta<1; all weighted A\_\allowbreak{}g,V\_\allowbreak{}g extend boundedly.}\par
{\footnotesize \textit{Remaining.} The normalized inverse, positive fast form and unique bounded Schur minimum are formalized. Identify actual Wilson form domains and sources;}\par
\smallskip
\noindent\textbf{\texttt{\small DERIV:\allowbreak{}WILSON\_\allowbreak{}SPATIAL\_\allowbreak{}SCHUR\_\allowbreak{}EXCESS:\allowbreak{}SP20\_\allowbreak{}SP24}} \hfill {\footnotesize conditional}\par
{\small Complete Schur comparisons give the reciprocal-gap recursion, product/inverse-fast-energy lower bound SP21-SP22 and source-overlap telescoping SP24 in common physical clock units.}\par
{\footnotesize \srcpath{docs/derivations/wilson-spatial-schur-excess.md} --- \#\# 6. Iterating energies and observable amplitudes in the physical clock (lines 412-477)}\par
{\footnotesize \textit{Hypotheses.} Closed self-adjoint reference and interacting forms on the stated compatible P/Q domains, already expressed in a common physical source chart and vacuum-subtracted.; Real energy z below the full fast restriction; positive retained energy weight B.; 0<alpha\_\allowbreak{}j<=1; full fast floors f\_\allowbreak{}j; summable epsilon\_\allowbreak{}j=1-alpha\_\allowbreak{}j bounded below 1; summable inverse fast floors; actual source maps and errors eta\_\allowbreak{}j.}\par
{\footnotesize \textit{Remaining.} Formalize iteration of the closed-form gap theorem, positive infinite products and complete carrier/source projection transport.}\par
\smallskip
\noindent\textbf{\texttt{\small DERIV:\allowbreak{}WILSON\_\allowbreak{}SPATIAL\_\allowbreak{}SCHUR\_\allowbreak{}EXCESS:\allowbreak{}SP7\_\allowbreak{}SP10}} \hfill {\footnotesize conditional}\par
{\small K1=J\_\allowbreak{}z*W1J\_\allowbreak{}z and K2=J\_\allowbreak{}z*W2J\_\allowbreak{}z-V1*V1 are the complete Schur coefficients; SP8 gives the explicit O(|g|\textasciicircum{}3) remainder SP9, including the dressed force correction SP10.}\par
{\footnotesize \srcpath{docs/derivations/wilson-spatial-schur-excess.md} --- \#\# 3. The complete second-order coefficient and a uniform remainder (lines 115-182)}\par
{\footnotesize \textit{Hypotheses.} Closed self-adjoint reference and interacting forms on the stated compatible P/Q domains, already expressed in a common physical source chart and vacuum-subtracted.; Real energy z below the full fast restriction; positive retained energy weight B.; SP8 remainder constants a3,v2,v1,c1,g0 are finite and c1|g|<=theta<1.}\par
{\footnotesize \textit{Remaining.} The genuine Neumann operator expansion and cubic norm/sandwich remainder are formalized. Assemble the g-dependent direct term, moving forces, fast defect and spatial coefficients on the actual compatible Wilson form domain.}\par
\smallskip
\noindent\textbf{\texttt{\small DERIV:\allowbreak{}WILSON\_\allowbreak{}TRUE\_\allowbreak{}VACUUM\_\allowbreak{}BLOCK\_\allowbreak{}ESTIMATES:\allowbreak{}BA9\_\allowbreak{}BA10}} \hfill {\footnotesize conditional}\par
{\small The true conditional operator includes outside pressure W\_\allowbreak{}B=(H\_\allowbreak{}out Omega)/Omega and has gap >=gap(h\_\allowbreak{}B(b))-osc\_\allowbreak{}B W\_\allowbreak{}B; a frozen-block gap alone omits this term.}\par
{\footnotesize \srcpath{docs/derivations/wilson-true-vacuum-block-estimates.md} --- \#\# BA4. What a frozen-block spectral computation would additionally require (lines 142-168)}\par
{\footnotesize \textit{Hypotheses.} Finite periodic cubic SU(2) lattice with nondegenerate plaquettes and side L >= 3.; Product unit-S3 electric metric; epsilon > 0, k >= 0, lambda = k/epsilon; actual positive normalized Wilson ground Omega.; All conditional blocks use the actual vacuum law mu=Omega\textasciicircum{}2dU, not isolated-block spectral data.; The frozen-block gap and finite pressure oscillation are available on compatible form domains.}\par
{\footnotesize \textit{Remaining.} Formalize the conditional ground transform and bounded-potential min-max comparison, including the domain of H\_\allowbreak{}out Omega.}\par
\smallskip
\noindent\textbf{\texttt{\small DERIV:\allowbreak{}WILSON\_\allowbreak{}VACUUM\_\allowbreak{}ALIGNED\_\allowbreak{}ASSEMBLY:\allowbreak{}VA12\_\allowbreak{}VA13}} \hfill {\footnotesize conditional}\par
{\small For physical conditional-angle row kappa<1, G\_\allowbreak{}mu\textasciicircum{}2>=(1-kappa)G\_\allowbreak{}mu and gap\_\allowbreak{}phys(H)>=gamma\_\allowbreak{}*(1-kappa); the single plaquette physical restriction has E\_\allowbreak{}e=E0 and zero pair angle.}\par
{\footnotesize \srcpath{docs/derivations/wilson-vacuum-aligned-assembly.md} --- \#\# VA5. Projection geometry, with the Gauss law retained (lines 211-251)}\par
{\footnotesize \textit{Hypotheses.} Finite periodic cubic SU(2) lattice with nondegenerate plaquettes and side L >= 3.; Product unit-S3 electric metric; epsilon > 0, k >= 0, lambda = k/epsilon; actual positive normalized Wilson ground Omega.; c\_\allowbreak{}ef\textasciicircum{}phys is the actual L2 conditional-projection norm, and sup\_\allowbreak{}e sum\_\allowbreak{}(f!=e)c\_\allowbreak{}ef\textasciicircum{}phys<=kappa<1.}\par
{\footnotesize \textit{Remaining.} The arbitrary-dimensional bounded projection-sum implication and kernel-complement gap are formalized. Construct actual conditional projections, derive the local anticommutator bound from their Friedrichs angles, identify the common kernel as constants and establish the required row budget below one.}\par
\smallskip
\noindent\textbf{\texttt{\small DERIV:\allowbreak{}WILSON\_\allowbreak{}VACUUM\_\allowbreak{}ALIGNED\_\allowbreak{}ASSEMBLY:\allowbreak{}VA14\_\allowbreak{}VA19}} \hfill {\footnotesize conditional}\par
{\small Approximate tensorization or a summable block-angle row gives the stated physical gap with covering multiplicity; the original VA10 local floor tends to zero on the displayed logarithmic continuum trajectory.}\par
{\footnotesize \srcpath{docs/derivations/wilson-vacuum-aligned-assembly.md} --- \#\# VA6. The next estimate and the attempts to establish it (lines 252-415)}\par
{\footnotesize \textit{Hypotheses.} Finite periodic cubic SU(2) lattice with nondegenerate plaquettes and side L >= 3.; Product unit-S3 electric metric; epsilon > 0, k >= 0, lambda = k/epsilon; actual positive normalized Wilson ground Omega.; Actual conditional blocks, gamma\_\allowbreak{}min>0 and finite covering multiplicity m\_\allowbreak{}*; a spatial envelope is an additional hypothesis, not assigned data.}\par
{\footnotesize \textit{Remaining.} Actual-measure covariance minorization and explicit-domain unitary gap transport are formalized. Establish the actual Wilson conditional density floor, conditional disintegration, covariance-to-projection-angle norm and spatial row-budget bounds.}\par
\smallskip
\noindent\textbf{\texttt{\small DERIV:\allowbreak{}WILSON\_\allowbreak{}WEIGHTED\_\allowbreak{}REPAIR\_\allowbreak{}AND\_\allowbreak{}ROTOR\_\allowbreak{}GAP:\allowbreak{}WR17\_\allowbreak{}WR21}} \hfill {\footnotesize conditional}\par
{\small The actual trial comparison needs gamma\_\allowbreak{}Lambda-(mean R-inf R)>0 uniformly; a gauge-invariant patch has full-boundary trial excitations of energy 3epsilon, so an unrestricted boundary fast floor of order sqrt(epsilon k) cannot be inferred.}\par
{\footnotesize \srcpath{docs/derivations/wilson-weighted-repair-and-rotor-gap.md} --- \#\# WR6. The precise remaining derivation (lines 229-308)}\par
{\footnotesize \textit{Hypotheses.} Each operator is identified separately: classical quotient diffusion, complete compact S3 quantum rotor, actual coupled link lattice or independent product model as stated in the section.; Positive normalized overlapping-plaquette trial; gamma\_\allowbreak{}Lambda is an actual weighted generator gap.}\par
{\footnotesize \textit{Remaining.} Formalize the common form domains and mean-residual comparison; construct Haar-link marginal trial vectors proving the full-boundary obstruction.}\par
\smallskip
\noindent\textbf{\texttt{\small DERIV:\allowbreak{}YANGMILLS\_\allowbreak{}CONTINUUM\_\allowbreak{}BALABAN\_\allowbreak{}MULTISCALE\_\allowbreak{}PROOF:\allowbreak{}GHOST\_\allowbreak{}5\_\allowbreak{}8}} \hfill {\footnotesize conditional}\par
{\small Rescaling z=g*zeta exposes g\textasciicircum{}(1-kappa\_\allowbreak{}prime) in (5.8), but the required volume/scale-uniform covariant inverse-gradient bound is not proved by that algebra.}\par
{\footnotesize \srcpath{docs/derivations/yangmills-continuum-balaban-multiscale-proof.md} --- \#\#\# 5.2 Faddeev--Popov Ghost Determinant (lines 198-208)}\par
{\footnotesize \textit{Hypotheses.} Submitted SU(N), N>=2 continuum multiscale construction; read together with its in-document Repair record audited 2026-09-09.; The source title and summary claim completion, but its repair record retains explicit averaging/interpolation, irrelevance, physical-scale and scattering obligations. No machine certification is inferred from prose.; 0<kappa\_\allowbreak{}prime<1/2; a||zeta||\_\allowbreak{}infinity<=g\textasciicircum{}-kappa\_\allowbreak{}prime; the first operator inequality in (5.8) is an additional input.}\par
{\footnotesize \textit{Remaining.} Construct the precise ghost operator, determinant class and uniform inverse-gradient estimate on its constrained domain; scalar rescaling alone does not discharge them.}\par
\smallskip
\noindent\textbf{\texttt{\small DERIV:\allowbreak{}YANGMILLS\_\allowbreak{}CONTINUUM\_\allowbreak{}BALABAN\_\allowbreak{}MULTISCALE\_\allowbreak{}PROOF:\allowbreak{}OS0}} \hfill {\footnotesize conditional}\par
{\small The asserted Schwinger growth bound (8.2) supplies OS regularity only after its actual uniform moment/Sobolev estimates and continuum fields are constructed.}\par
{\footnotesize \srcpath{docs/derivations/yangmills-continuum-balaban-multiscale-proof.md} --- \#\#\# 8.1 OS0: Analyticity and Temperedness (lines 368-374)}\par
{\footnotesize \textit{Hypotheses.} Submitted SU(N), N>=2 continuum multiscale construction; read together with its in-document Repair record audited 2026-09-09.; The source title and summary claim completion, but its repair record retains explicit averaging/interpolation, irrelevance, physical-scale and scattering obligations. No machine certification is inferred from prose.}\par
{\footnotesize \textit{Remaining.} Formalize the continuum Schwinger distributions and prove the factorial/test-seminorm bound from the actual regulated measures.}\par
\smallskip
\noindent\textbf{\texttt{\small DERIV:\allowbreak{}YANGMILLS\_\allowbreak{}CONTINUUM\_\allowbreak{}BALABAN\_\allowbreak{}MULTISCALE\_\allowbreak{}PROOF:\allowbreak{}OS3}} \hfill {\footnotesize conditional}\par
{\small Bosonic multiplication observables have permutation-symmetric Schwinger functions, provided the claimed continuum distributions exist.}\par
{\footnotesize \srcpath{docs/derivations/yangmills-continuum-balaban-multiscale-proof.md} --- \#\#\# 8.4 OS3: Permutation Symmetry (lines 409-411)}\par
{\footnotesize \textit{Hypotheses.} Submitted SU(N), N>=2 continuum multiscale construction; read together with its in-document Repair record audited 2026-09-09.; The source title and summary claim completion, but its repair record retains explicit averaging/interpolation, irrelevance, physical-scale and scattering obligations. No machine certification is inferred from prose.}\par
{\footnotesize \textit{Remaining.} Formalize products of the constructed fields/observables and the permutation action on their distributional arguments.}\par
\smallskip
\noindent\textbf{\texttt{\small DERIV:\allowbreak{}YANGMILLS\_\allowbreak{}CONTINUUM\_\allowbreak{}BALABAN\_\allowbreak{}MULTISCALE\_\allowbreak{}PROOF:\allowbreak{}OS\_\allowbreak{}RECONSTRUCTION}} \hfill {\footnotesize conditional}\par
{\small With corrected OS hypotheses and totality, the reflection-positive quotient completion has a physical Hilbert space, vacuum and nonnegative self-adjoint time generator.}\par
{\footnotesize \srcpath{docs/derivations/yangmills-continuum-balaban-multiscale-proof.md} --- \#\#\# 9.1 The Physical Hilbert Space and Hamiltonian (lines 439-453)}\par
{\footnotesize \textit{Hypotheses.} Submitted SU(N), N>=2 continuum multiscale construction; read together with its in-document Repair record audited 2026-09-09.; The source title and summary claim completion, but its repair record retains explicit averaging/interpolation, irrelevance, physical-scale and scattering obligations. No machine certification is inferred from prose.}\par
{\footnotesize \textit{Remaining.} Formalize the OS quotient, completion, translations and corrected reconstruction regularity; actual OS inputs remain separate obligations.}\par
\smallskip
\noindent\textbf{\texttt{\small DERIV:\allowbreak{}YANGMILLS\_\allowbreak{}CONTINUUM\_\allowbreak{}BALABAN\_\allowbreak{}MULTISCALE\_\allowbreak{}PROOF:\allowbreak{}SCATTERING}} \hfill {\footnotesize conditional}\par
{\small Haag-Ruelle scattering is conditional on an isolated positive single-particle mass shell; the source now explicitly states that dispersion hypothesis. Nonzero off-shell/tree-level four-point data alone do not supply the required on-shell amplitude lower bound.}\par
{\footnotesize \srcpath{docs/derivations/yangmills-continuum-balaban-multiscale-proof.md} --- \#\#\# 9.3 Relativistic Wightman Fields and Non-Trivial Scattering (\$S \textbackslash{}ne I\$) (lines 479-501)}\par
{\footnotesize \textit{Hypotheses.} Submitted SU(N), N>=2 continuum multiscale construction; read together with its in-document Repair record audited 2026-09-09.; The source title and summary claim completion, but its repair record retains explicit averaging/interpolation, irrelevance, physical-scale and scattering obligations. No machine certification is inferred from prose.; Isolated mass hyperboloid separated below multiparticle threshold; locality, spectrum and actual LSZ/Haag-Ruelle assumptions.}\par
{\footnotesize \textit{Remaining.} Formalize the isolated-shell hypothesis, asymptotic states and a controlled on-shell nonzero amplitude; establish these for the constructed theory before concluding S!=I.}\par
\smallskip
\noindent\textbf{\texttt{\small DERIV:\allowbreak{}YANGMILLS\_\allowbreak{}CONTINUUM\_\allowbreak{}BALABAN\_\allowbreak{}MULTISCALE\_\allowbreak{}PROOF:\allowbreak{}THEOREM\_\allowbreak{}4\_\allowbreak{}2}} \hfill {\footnotesize conditional}\par
{\small The large-field damping claim (4.4)-(4.5) is conditional on a positive coarse-action subtraction constant and the unproved irrelevance gain identified in the repair record.}\par
{\footnotesize \srcpath{docs/derivations/yangmills-continuum-balaban-multiscale-proof.md} --- \#\#\# 4.2 Super-Exponential Suppression of Large Fields (lines 157-180)}\par
{\footnotesize \textit{Hypotheses.} Submitted SU(N), N>=2 continuum multiscale construction; read together with its in-document Repair record audited 2026-09-09.; The source title and summary claim completion, but its repair record retains explicit averaging/interpolation, irrelevance, physical-scale and scattering obligations. No machine certification is inferred from prose.; 0<kappa\_\allowbreak{}prime<kappa<1/6; actual regularized functional measure and corrected background construction; uniform constants.}\par
{\footnotesize \textit{Remaining.} Prove the quantitative irrelevance/subtraction constant and the actual large-field measure estimate, including the averaging constraint and fluctuation direction invisible to a coarse penalty.}\par
\smallskip
\noindent\textbf{\texttt{\small DERIV:\allowbreak{}YANGMILLS\_\allowbreak{}CONTINUUM\_\allowbreak{}BALABAN\_\allowbreak{}MULTISCALE\_\allowbreak{}PROOF:\allowbreak{}THEOREM\_\allowbreak{}6\_\allowbreak{}2}} \hfill {\footnotesize conditional}\par
{\small The claimed RG contraction (6.5) has the repaired induction envelope ||K\_\allowbreak{}k||<=4 C\_\allowbreak{}ind g\_\allowbreak{}k\textasciicircum{}2, but its single-block contraction gain still requires the irrelevance lemma.}\par
{\footnotesize \srcpath{docs/derivations/yangmills-continuum-balaban-multiscale-proof.md} --- \#\#\# 6.3 The Mayer Tree-Graph Expansion \& Cluster Contraction (lines 269-295)}\par
{\footnotesize \textit{Hypotheses.} Submitted SU(N), N>=2 continuum multiscale construction; read together with its in-document Repair record audited 2026-09-09.; The source title and summary claim completion, but its repair record retains explicit averaging/interpolation, irrelevance, physical-scale and scattering obligations. No machine certification is inferred from prose.; lambda\_\allowbreak{}contraction in (0,1/2); sufficiently small initial g0; the corrected averaging, fluctuation and large-field inputs hold.}\par
{\footnotesize \textit{Remaining.} Formalize the actual tree expansion and prove the residual local-operator irrelevance gain after marginal subtraction; then prove the model-dependent contraction.}\par
\smallskip
\noindent\textbf{\texttt{\small DERIV:\allowbreak{}YANGMILLS\_\allowbreak{}CONTINUUM\_\allowbreak{}BALABAN\_\allowbreak{}MULTISCALE\_\allowbreak{}PROOF:\allowbreak{}THEOREM\_\allowbreak{}6\_\allowbreak{}3}} \hfill {\footnotesize conditional}\par
{\small The source records b0=11N/(48pi\textasciicircum{}2), the perturbative running formula (6.11)-(6.12) and a uniform polymer conclusion (Corollary 6.4); its repaired induction requires quantitative remainder control and a propagated envelope.}\par
{\footnotesize \srcpath{docs/derivations/yangmills-continuum-balaban-multiscale-proof.md} --- \#\#\# 6.4 Asymptotic Freedom and Infinite Scale Stability (lines 296-314)}\par
{\footnotesize \textit{Hypotheses.} Submitted SU(N), N>=2 continuum multiscale construction; read together with its in-document Repair record audited 2026-09-09.; The source title and summary claim completion, but its repair record retains explicit averaging/interpolation, irrelevance, physical-scale and scattering obligations. No machine certification is inferred from prose.; Perturbative UV regime; controlled beta remainder and actual RG map; this does not place the hadronic scale in the small-g domain.}\par
{\footnotesize \textit{Remaining.} Formalize the actual RG beta coefficient and remainder assumptions before asserting monotonicity for all g0; prove the stated envelope with the repaired constants.}\par
\smallskip
\noindent\textbf{\texttt{\small DERIV:\allowbreak{}YANGMILLS\_\allowbreak{}CONTINUUM\_\allowbreak{}BALABAN\_\allowbreak{}MULTISCALE\_\allowbreak{}PROOF:\allowbreak{}THEOREM\_\allowbreak{}7\_\allowbreak{}2}} \hfill {\footnotesize conditional}\par
{\small The continuum measure claim requires a continuous positive-definite characteristic functional on the full nuclear test space; convergence only of gauge-invariant observables is not itself this construction.}\par
{\footnotesize \srcpath{docs/derivations/yangmills-continuum-balaban-multiscale-proof.md} --- \#\#\# 7.2 Construction of the Limiting Continuum Measure \$d\textbackslash{}mu\_\allowbreak{}\{\textbackslash{}text\{YM\}\}\$ (lines 341-363)}\par
{\footnotesize \textit{Hypotheses.} Submitted SU(N), N>=2 continuum multiscale construction; read together with its in-document Repair record audited 2026-09-09.; The source title and summary claim completion, but its repair record retains explicit averaging/interpolation, irrelevance, physical-scale and scattering obligations. No machine certification is inferred from prose.; Theorem 7.1 with the repaired summable bound; actual uniform regularity and characteristic-functional estimates.}\par
{\footnotesize \textit{Remaining.} Construct the characteristic functional with continuity/positivity and a consistent gauge-invariant observable realization, then apply Minlos-Bochner.}\par
\smallskip
\noindent\textbf{\texttt{\small DERIV:\allowbreak{}YANGMILLS\_\allowbreak{}CONTINUUM\_\allowbreak{}BALABAN\_\allowbreak{}MULTISCALE\_\allowbreak{}PROOF:\allowbreak{}THEOREM\_\allowbreak{}8\_\allowbreak{}1}} \hfill {\footnotesize conditional}\par
{\small Centered reflection-adapted deterministic blocking preserves reflection positivity under its exact support/equivariance hypotheses, and a convergent compatible limit retains the inequality (8.5).}\par
{\footnotesize \srcpath{docs/derivations/yangmills-continuum-balaban-multiscale-proof.md} --- \#\#\# 8.3 OS2: Reflection Positivity (lines 383-408)}\par
{\footnotesize \textit{Hypotheses.} Submitted SU(N), N>=2 continuum multiscale construction; read together with its in-document Repair record audited 2026-09-09.; The source title and summary claim completion, but its repair record retains explicit averaging/interpolation, irrelevance, physical-scale and scattering obligations. No machine certification is inferred from prose.; Reflection-adapted deterministic block map, positive-time support and actual convergent correlations; the cited centered-blocking construction applies.}\par
{\footnotesize \textit{Remaining.} Formalize the deterministic reflected blocking and limit passage for the actual continuum measure; general reflection-equivariant Markov blocking is not this theorem.}\par
\smallskip
\noindent\textbf{\texttt{\small DERIV:\allowbreak{}YANGMILLS\_\allowbreak{}CONTINUUM\_\allowbreak{}BALABAN\_\allowbreak{}MULTISCALE\_\allowbreak{}PROOF:\allowbreak{}THEOREM\_\allowbreak{}9\_\allowbreak{}1}} \hfill {\footnotesize conditional}\par
{\small A common strictly positive exponential physical-time correlation rate on the dense centered reconstructed family implies spectrum(H) subset \{0\} union [Delta\_\allowbreak{}phys,infinity), with the simple vacuum; this conditional spectral argument is confirmed sound by the source repair record.}\par
{\footnotesize \srcpath{docs/derivations/yangmills-continuum-balaban-multiscale-proof.md} --- \#\#\# 9.2 The Physical Spectral Mass Gap (lines 454-478)}\par
{\footnotesize \textit{Hypotheses.} Submitted SU(N), N>=2 continuum multiscale construction; read together with its in-document Repair record audited 2026-09-09.; The source title and summary claim completion, but its repair record retains explicit averaging/interpolation, irrelevance, physical-scale and scattering obligations. No machine certification is inferred from prose.; Actual reconstructed physical space with dense observable family; common rate Delta\_\allowbreak{}phys>0 in physical time; simple vacuum.}\par
{\footnotesize \textit{Remaining.} Formalize the spectral-measure positivity, infinite-time exclusion and density extension; derive the actual physical clustering input (9.7) separately.}\par
\smallskip
\noindent\textbf{\texttt{\small DERIV:\allowbreak{}YANGMILLS\_\allowbreak{}GPU\_\allowbreak{}RESOLUTION\_\allowbreak{}AUDIT:\allowbreak{}GA8\_\allowbreak{}GA9}} \hfill {\footnotesize conditional}\par
{\small The inserted scalar recursion has explicit positive product/tail bounds GA8 when its assumptions hold; the Phase 4 four-plaquette builder remains a tensor sum GA9.}\par
{\footnotesize \srcpath{docs/derivations/yangmills-gpu-resolution-audit.md} --- \#\# GA4. The additional Phase 4 has a valid conditional scalar limit (lines 162-209)}\par
{\footnotesize \textit{Hypotheses.} The preserved September 9 scripts are inspected as the operators they actually implement; they are not assumed to represent Wilson dynamics.; Analytic matrix identities, replayed floating observations and reported-only runs remain distinct.; eps\_\allowbreak{}j=C(ell+jh)\textasciicircum{}(-3/2)<1; positive initial gap and fast floors; these are supplied inputs in the script.}\par
{\footnotesize \textit{Remaining.} Formalize the scalar positive-product and geometric tail bounds; derive actual coarse Wilson parameters before applying them.}\par
\smallskip
\noindent\textbf{\texttt{\small DERIV:\allowbreak{}YANGMILLS\_\allowbreak{}RECONSTRUCTION:\allowbreak{}R6\_\allowbreak{}R7}} \hfill {\footnotesize conditional}\par
{\small If C\_\allowbreak{}F(t)<=A\_\allowbreak{}F exp(-eta*t+alpha*R)+B\_\allowbreak{}F exp(-beta*R) for every R,t>=0, choose R=eta*t/(alpha+beta) to obtain common rate eta*beta/(alpha+beta).}\par
{\footnotesize \srcpath{docs/derivations/yangmills-reconstruction.md} --- \#\# 5. A precise sufficient repair, with its remaining hypothesis exposed (lines 90-111)}\par
{\footnotesize \textit{Hypotheses.} A nonnegative self-adjoint physical Hamiltonian with a specified vacuum and its positive spectral measures; physical Euclidean time is used.; OS reconstruction and any limit theorem are explicit analytic inputs, not consequences of the finite controls.; eta,beta>0,alpha>=0; one common rate and controlled cutoff/volume constants on a dense centered physical family.}\par
{\footnotesize \textit{Remaining.} Real-radius exponent balancing, actual-measure spectral exclusion and dense-family extension are formalized. Establish the actual uniform Yang-Mills two-parameter localization estimate, physical spectral identification and total family;}\par
\smallskip
\noindent\textbf{\texttt{\small DERIV:\allowbreak{}YANGMILLS\_\allowbreak{}RESOLVENT\_\allowbreak{}LOCALIZATION\_\allowbreak{}AND\_\allowbreak{}DIRICHLET\_\allowbreak{}GAP:\allowbreak{}LOCAL\_\allowbreak{}KERNEL\_\allowbreak{}IMPLICATION}} \hfill {\footnotesize conditional}\par
{\small Given uniform massive exponentially decaying kernels and a finite-range perturbation of size C2|g|, the localized scalar pairing has a volume-independent geometric-sum bound.}\par
{\footnotesize \srcpath{docs/derivations/yangmills-resolvent-localization-and-dirichlet-gap.md} --- \#\#\# 2.3 Volume-Uniformity of the Pairing (lines 62-80)}\par
{\footnotesize \textit{Hypotheses.} This submitted proposal is explicitly preserved under its September 9 source audit; closure claims are not the current graph verdict.; The actual model/operator identification is an obligation, separate from a scalar finite replay.; Kernel decay constants and perturbation range/row bounds uniform in volume; finite localized source with the source normalization explicitly retained.}\par
{\footnotesize \textit{Remaining.} Infinite kernel-pairing summability, exact integer-lattice sums and the full source l1 factor are formalized. Identify actual Wilson fast kernels and connected signed sources and prove the required uniform decay and perturbation envelopes.}\par
\smallskip
\noindent\textbf{\texttt{\small DERIV:\allowbreak{}YANGMILLS\_\allowbreak{}RESOLVENT\_\allowbreak{}LOCALIZATION\_\allowbreak{}AND\_\allowbreak{}DIRICHLET\_\allowbreak{}GAP:\allowbreak{}RESOLVENT\_\allowbreak{}IDENTITY}} \hfill {\footnotesize conditional}\par
{\small Compatible invertible fast operators obey R\_\allowbreak{}g-R0=-R0(F\_\allowbreak{}g-F0)R\_\allowbreak{}g and the stated selected-pairing identity.}\par
{\footnotesize \srcpath{docs/derivations/yangmills-resolvent-localization-and-dirichlet-gap.md} --- \#\#\# 2.1 The Resolvent Identity (lines 42-54)}\par
{\footnotesize \textit{Hypotheses.} This submitted proposal is explicitly preserved under its September 9 source audit; closure claims are not the current graph verdict.; The actual model/operator identification is an obligation, separate from a scalar finite replay.; Real z below both spectra; compatible domains and adjoint pairings.}\par
{\footnotesize \textit{Remaining.} Formalize the resolvent identity on the compatible closed-form domains.}\par
\smallskip
\noindent\textbf{\texttt{\small DERIV:\allowbreak{}YANGMILLS\_\allowbreak{}WEIGHTED\_\allowbreak{}CURVATURE:\allowbreak{}GST3}} \hfill {\footnotesize conditional}\par
{\small Ric+(2/hbar)Hess S>=kappa gives ||A\_\allowbreak{}S f||\textasciicircum{}2>=c*kappa<f,A\_\allowbreak{}Sf>, hence spectrum in \{0\} union [c*kappa,infinity) and the Poincare bound when the kernel is constants.}\par
{\footnotesize \srcpath{docs/derivations/yangmills-weighted-curvature.md} --- \#\# GST-3: curvature normalization and a spectral argument without a discrete excited state (lines 92-130)}\par
{\footnotesize \textit{Hypotheses.} Finite-dimensional smooth Riemannian configuration space or the explicitly specified finite matrix/oscillator model.; H=-c Delta+V, c=hbar\textasciicircum{}2/(2m)>0; psi=N exp(-S/hbar)>0 normalized; common operator/form domains as stated.; kappa>0; closed Bochner identity/inequality and simple constant kernel.}\par
{\footnotesize \textit{Remaining.} The bounded positive-operator square-to-gap implication is formalized in arbitrary complex Hilbert dimension. Construct the actual weighted operator, verify its quadratic inequality and manage the unbounded-domain/spectral realization required by the source.}\par
\smallskip
\noindent\textbf{\texttt{\small DERIV:\allowbreak{}YANGMILLS\_\allowbreak{}WEIGHTED\_\allowbreak{}CURVATURE:\allowbreak{}GST6}} \hfill {\footnotesize conditional}\par
{\small The closed-form min-max comparison yields E1(H)-E0(H)>=gamma-(mean R-rmin).}\par
{\footnotesize \srcpath{docs/derivations/yangmills-weighted-curvature.md} --- \#\# GST-6: the residual-mean certificate (lines 167-190)}\par
{\footnotesize \textit{Hypotheses.} Finite-dimensional smooth Riemannian configuration space or the explicitly specified finite matrix/oscillator model.; H=-c Delta+V, c=hbar\textasciicircum{}2/(2m)>0; psi=N exp(-S/hbar)>0 normalized; common operator/form domains as stated.; A\_\allowbreak{}S is nonnegative self-adjoint with simple normalized constant ground and first variational level >=gamma>0; R>=rmin, finite mean; 1 belongs to the sum-form domain; min-max applies.}\par
{\footnotesize \textit{Remaining.} Formalize the lower first-excited variational comparison and the trial-vector Rayleigh upper bound on common form domains.}\par
\smallskip
\endgroup

